This chapter presents the research methodology employed for the development of a machine-learning-driven fuel monitoring and anomaly detection system for fleet vehicles. The methodology integrates a multi-layered approach combining OBD-II data acquisition, edge computing, IoT communication, cloud-based machine learning, and client-side visualization. The approach emphasizes real-time anomaly detection, accurate fuel monitoring, operational transparency, and actionable driver and fleet management insights.

The research methodology is structured to address both general and specific objectives, ensuring that the system is technically feasible, scalable, and reliable. Figure~\ref{fig:End-to-end_system_architecture} and Table~\ref{tab:methods_per_objective} summarize the methods and approaches for each objective.

The chapter placement was reviewed to keep the paper aligned with the DLSU thesis structure. Chapter~\ref{ch:theorycon} is retained for the theoretical basis of IoT telemetry, anomaly detection, and fleet monitoring; Chapter~\ref{ch:designcon} is retained for design decisions, standards, limits, and constraints; and this chapter focuses on the actual research procedure: how the prototype was implemented, how telemetry was collected and processed, how the machine-learning variants were trained, and how each objective was evaluated.
\begin{center}
\scriptsize
\begin{tabularx}{\textwidth}{p{0.2\textwidth}|X|p{0.1\textwidth}}
\caption{Summary of methods for reaching the objectives} \label{tab:methods_per_objective} \\
\hline 
\hline 
\textbf{Objectives} & \textbf{Methods / Approaches} & \textbf{Locations} \\ 
\hline 
GO: To develop and evaluate an IoT-based system for trucks, in collaboration with the partner truck company Hino Parañaque Subic GS Auto Inc. (SGSA), that monitors real-time fuel levels, tracks routes using GPS, and detects unusual fuel consumption using machine learning, with the goal of improving driver safety, supporting rest scheduling, and facilitating proactive vehicle maintenance through automated alerts using fuel and distance metrics, based on established service guidelines. &
\begin{itemize}
    \item Embedded MCU preprocessing and aggregation
    \item LoRa/Cellular communication for low-latency and fail-safe data transfer
    \item Cloud-based machine learning for anomaly detection
    \item Web dashboard for visualization
    \item Driver alert subsystem for operational safety
\end{itemize} & Sec.~\ref{sec:implement} \\
\hline 
SO1: To measure fuel levels using digital float sensors integrated via the truck's ECU/diagnostic connector and GPS module, ensuring non-invasive installation and real-time accuracy within $\pm$5\%. &
\begin{itemize}
    \item Digital fuel-level acquisition through the truck's ECU/diagnostic connector
    \item Edge-side median, EMA, dead-band, and direction-aware filtering
    \item Validation against dashboard-gauge walkaround pairs and signal variance
    \item Continuous validation for measurement reliability
\end{itemize} & Secs.~\ref{sec:fuel_acquisition}, \ref{sec:fuel_preprocessing} \\
\hline 
SO2: To transmit real-time fuel and GPS data through IoT-based communication to a cloud server, accessible through a centralized dashboard. &
\begin{itemize}
    \item Real-time IoT transmission via LoRa, LTE, GSM 2G, WiFi, and SPIFFS replay
    \item Cloud ingestion APIs
    \item Synchronized HMI, mobile driver application, and administrative dashboard
    \item Latency measurement and benchmarking to ensure timely reporting
\end{itemize} & Sec.~\ref{sec:communication} \\
\hline 
SO3: To capture and log GPS location data every 5 seconds via an embedded module, correlating route information with fuel and time data for behavior tracking. &
\begin{itemize}
    \item GPS-based location logging at 5-second intervals
    \item Synchronization with fuel data
    \item Structured cloud database for multi-source data storage and route–fuel correlation
\end{itemize} & Secs.~\ref{sec:implement}, \ref{sec:backend} \\
\hline 
SO4: To design and evaluate a context-aware fuel usage deviation model using non-parametric unsupervised machine learning models (Matrix Profile or Isolation Forest), achieving 90\% precision and 10\% false positives in identifying anomalous fuel consumption patterns within controlled environments. &
\begin{itemize}
    \item Offline training of anomaly detection models using Isolation Forest and Matrix Profile
    \item Evaluation using precision, FPR, and accuracy
    \item Context-aware thresholding to reduce false positives
\end{itemize} & Secs.~\ref{sec:ml}, \ref{sec:evaluation} \\
\hline 
SO5: To develop and integrate hardware-based rest and maintenance alert mechanisms (e.g., Start/Stop button, buzzer indicators) tied to the web dashboard using cumulative fuel and distance data. &
\begin{itemize}
    \item Driver rest and maintenance alert system with cumulative distance/time tracking
    \item Cloud-based monitoring and logging
    \item Alert notification via dashboard and in-vehicle signals
\end{itemize} & Secs.~\ref{sec:hardware}, \ref{sec:alerts} \\
\hline
\end{tabularx}
\end{center}

Table~\ref{tab:methods_per_objective} is used as the methodological roadmap for the chapter. Each row connects a thesis objective to the concrete subsystem, data source, or validation activity used to satisfy it. The table also prevents the methodology from becoming only an implementation narrative: it shows that fuel acquisition, communication, GPS logging, machine learning, and alert delivery are each tied to a measurable requirement that is later evaluated in Chapter~\ref{ch:result_discuss}.


\section{System Architecture and Conceptual Design}
\label{sec:implement}

\begin{figure}[!htbp]
	\centering
	\includegraphics[width=1\textwidth]{concep_diagram.png}
	\caption{Conceptual Framework}
	\label{fig:conceptual_framework}
\end{figure}

The system follows a layered architecture with a focus on modularity, scalability, and reliability. Figure~\ref{fig:conceptual_framework} integrates the following:

\begin{itemize}
    \item \textbf{Vehicle Sensing Subsystem:} OBD-II fuel and vehicle-context acquisition, embedded GPS, mobile GPS fallback, and driver rest-time monitors.
    \item \textbf{Edge Processing:} The MCU performs signal preprocessing, filtering, smoothing, synchronization, and timestamping before transmission.
    \item \textbf{IoT Communication:} LoRa, LTE, GSM 2G, WiFi, and SPIFFS replay provide prioritized telemetry delivery under changing network conditions.
    \item \textbf{Cloud Intelligence:} Isolation Forest Model C and Matrix Profile evaluate incoming telemetry and trigger fuel-anomaly alerts when the contextual gate permits.
    \item \textbf{Client Visualization:} The embedded HMI, mobile driver web application, and administrative dashboard provide synchronized fuel, route, trip-state, and alert views.
\end{itemize}

The purpose of the conceptual framework is to show the boundary of the design solution before individual subsystems are discussed. The vehicle layer produces the measurements, the edge layer cleans and packages them, the communication layer preserves delivery under changing coverage, the cloud layer stores and interprets the data, and the interface layer returns the interpreted state to the driver and fleet manager. This layered reading is important because the prototype is evaluated as an end-to-end monitoring system rather than as a single sensor or dashboard.

The system architecture in Figure~\ref{fig:End-to-end_system_architecture} is an edge-to-cloud fleet-monitoring pipeline in which the embedded HMI performs data acquisition, filtering, trip-state handling, local driver display, and prioritised transmission. The backend stores telemetry, synchronises state across three concurrent client surfaces - the in-vehicle HMI, the mobile driver web application, and the administrative dashboard - and dispatches anomaly and safety alerts through the same database tables that the three surfaces subscribe to.

\FloatBarrier
\begin{figure}[!htbp]
    \centering
        \includegraphics[width=0.8\textwidth]{End-to-end_system_architecture.png}
    \caption{End-to-end System Architecture}
    \label{fig:End-to-end_system_architecture}
\end{figure}
\clearpage

Figure~\ref{fig:End-to-end_system_architecture} illustrates the end-to-end communication architecture of the system, depicting how data flows from the ESP32 HMI through multiple transmission channels to the cloud backend and connected client interfaces.  The HMI, serving as edge node, collects fuel, location, and alert data while also providing the user interface for the driver.  From the HMI, data is transmitted through WiFi, LTE/GSM 2G, and LoRa before converging at the Supabase and Render backend, which handles telemetry ingestion, alert dispatching, and ML anomaly detection. The SPIFFS onboard flash buffer serves as a local fallback wherein stored telemetry is synced when a connection is re-established. The fallback priority order (LoRa $\rightarrow$ LTE $\rightarrow$ GSM 2G $\rightarrow$ WiFi $\rightarrow$ SPIFFS buffer $\rightarrow$ delayed sync) ensures continuous data delivery even under degraded network conditions. From the backend, processed data is propagated to the mobile driver web application, the administrative dashboard, and the HMI poll-back, maintaining real-time synchronization across all three client surfaces.

The architecture also identifies where each later measurement is taken. Fuel accuracy is measured before and after the edge filter, communication latency is measured between device emission and backend receipt, GPS cadence is measured at the device timestamp, anomaly detection is measured after backend feature construction, and alert usability is measured at the HMI, mobile, and dashboard surfaces. This placement makes the Chapter~6 result chain traceable from sensor input to user-facing output.

\subsection{Input-Process-Output (IPO) Overview}

The data design outlines a cohesive pipeline that integrates OBD-II acquisition, GPS synchronization, cloud transmission, and model-based analysis to enable real-time anomaly detection in fleet fuel systems. Data collection begins with the embedded HMI capturing OBD-II fuel level, vehicle context, GPS coordinates, communication-channel metadata, and trip-state events. These signals are preprocessed through the fuel filter cascade and synchronized before backend ingestion.

To ensure temporal alignment, telemetry records include device-side and backend-side timestamps. This synchronization step is critical for comparing trends across OBD-II, GPS, communication, and alert data. The processed signals are uploaded to the Supabase and Render backend for storage, anomaly scoring, and synchronized access across the HMI, mobile driver application, and administrative dashboard.

For anomaly detection, the synchronized time-series data are used to train and evaluate Isolation Forest and Matrix Profile detectors. These models are tested on project-native trips and controlled-injection scenarios to reflect normal and anomalous behavior. Key performance metrics, including MAPE\glsadd{MAPE}, timestamp drift, detection latency, precision, recall, and false positive rate, are used to assess system effectiveness.

\subsubsection{IPO Chart}

\begin{figure}[!htbp]
    \centering
        \includegraphics[width=1\textwidth]{IPO.png}
    \caption{IPO Chart}
    \label{fig:ipo_chart}
\end{figure}

The IPO (Input-Process-Output) chart in Figure~\ref{fig:ipo_chart} provides a structured overview of the IoT-based fleet fuel monitoring and anomaly detection system. The input data consists of real-time readings from the fleet vehicle's CAN bus/OBD-II fuel signal, GPS module coordinates and speed at five-second telemetry intervals, ignition status, and in-cab rest alert subsystem signals. These inputs are supplemented by historical fuel anomaly datasets and controlled-injection anomalies to simulate rare or edge-case events. Moderating variables such as driving behavior, network conditions, and environmental factors also influence the data.

The process begins with signal preprocessing, where the MCU filters and timestamps fuel, GPS, and ignition data to reduce noise and ensure synchronization. The preprocessed data is then transmitted through the prioritized communication stack and stored in the cloud backend. Anomaly detection is carried out in the cloud using Matrix Profile for time-series discord identification and Isolation Forest for feature-based outlier detection. Model validation and evaluation are performed using labelled and test datasets, calculating performance metrics such as precision, recall, false positive rate (FPR), detection latency, and overall accuracy.

The system outputs include fuel anomaly flags, such as potential theft or leakage, and driver rest alerts based on cumulative distance or driving duration. Additionally, it provides real-time GPS tracking on a web dashboard, estimated fuel levels as a percentage of tank capacity, and quantitative performance metrics for operational evaluation. This end-to-end flow ensures accurate, real-time monitoring, proactive safety notifications, and actionable insights for fleet management.

\subsection{End-to-End Telemetry Data Flow}

\begin{lstlisting}[caption=End-to-End Telemetry Data Flow,label=lst:telemetry_data_flow]
START
WHILE system is running DO
    READ ECU fuel data
    READ ECU operational data
    READ GPS data
    APPLY Median Filter
    APPLY EMA Filter
    APPLY Dead-Band Filter
    APPLY Direction-Aware Filter
    // UPDATE trip state
        IF vehicle is loading THEN
            trip_state = "Loading"
        ELSE IF vehicle is moving THEN
            trip_state = "Active"
        ELSE IF temporary stop detected THEN
            trip_state = "Paused"
        ELSE IF trip completed THEN
            trip_state = "Ended"
        END IF
    COMPUTE fused distance using GPS and OBD data
    // SELECT available communication channel
        IF LoRa available THEN
            Send via LoRa
        ELSE IF LTE available THEN
            Send via LTE
        ELSE IF GSM available THEN
            Send via GSM
        ELSE IF WiFi available THEN
            Send via WiFi
        ELSE
            Save data to SPIFFS
        END IF
    SEND telemetry to backend
    INSERT telemetry into telemetry_logs
    RUN ML anomaly detection
    IF anomaly passes contextual gate THEN
        INSERT alert into alerts
    END IF
    UPDATE database
    UPDATE dashboard, mobile app, and HMI display
END WHILE
END
\end{lstlisting}

Listing~\ref{lst:telemetry_data_flow} converts the architecture into the operating sequence followed by the prototype during a trip. It shows that telemetry collection is cyclical: the device reads fuel and GPS, filters the fuel signal, updates trip state, selects a communication channel, stores the record, runs anomaly detection, and then refreshes all user interfaces. The listing is not intended to expose source code; it is a readable process model showing how the hardware, backend, and interface layers cooperate during normal and degraded-network operation.

\section{Hardware Methodology}
\label{sec:hardware}

The embedded hardware platform is built around an ESP32 dual-core microcontroller running at 240 MHz with separate VSPI and HSPI\glsadd{HSPI} bus assignments to prevent SPI\glsadd{SPI} arbitration conflicts between the TFT display and the LoRa transceiver. Peripheral components include the ILI9341 2.4-inch TFT display with XPT2046 capacitive touch controller, physical action buttons, and a piezo buzzer. The MCP2515 CAN controller communicates with the ESP32 through SPI and interfaces with the vehicle CAN lines through a CAN transceiver.
The A7670E LTE/GSM modem and embedded GPS module use UART links, while the SX1278 LoRa transceiver uses VSPI.

\subsection{HMI Control Pins}

The main driver-facing controls use two debounced physical buttons for rest and snooze actions, one PWM\glsadd{PWM}-driven buzzer output for alerts, and one modem power-control pin. The full GPIO map is provided in Appendix~\ref{ch:appendix_technical_details} so the main methodology can focus on the hardware role rather than pin-level implementation detail.

This pin assignment supports the SO5 alert workflow by giving the driver physical controls that remain usable even when the touchscreen is difficult to operate during field conditions. The rest and snooze buttons are treated as event sources, the buzzer as the immediate alert actuator, and the modem-control line as the recovery mechanism for cellular communication faults.

\subsection{Hardware Block Diagram}

\begin{figure}[!htbp]
    \centering
    \includegraphics[width=0.9\textwidth]{image.png}
    \caption{Embedded HMI Hardware Block Diagram}
    \label{fig:esp32_hardware_block_diagram}
\end{figure}

Figure~\ref{fig:esp32_hardware_block_diagram} shows the hardware block diagram of the embedded HMI node, depicting how the ESP32 microcontroller connects to the OBD-II CAN interface, the A7670E cellular modem, the SX1278 LoRa transceiver, the GPS module, the ILI9341 TFT display, and the driver alert hardware on its respective bus and UART lines.

The block diagram should be read as the physical implementation counterpart of the system architecture. It explains why separate buses were used: the TFT display, CAN controller, LoRa transceiver, modem, and GPS receiver all generate time-sensitive events, so the wiring design separates high-traffic display updates from telemetry and radio activity. This separation reduces the chance that display refresh, CAN polling, and radio transmission interfere with one another during active trips.

\section{OBD-II / ECU Fuel and Vehicle Data Acquisition Methodology}
\label{sec:fuel_acquisition}

The HMI acquired vehicle telemetry through the standard OBD-II interface, which provides universal Mode 01 access to contextual variables --- vehicle speed, engine RPM, engine load, throttle position, mass air flow, coolant temperature, engine runtime, and diagnostic-trouble-code status --- on every OBD-II-compliant vehicle regardless of manufacturer. Fuel level is treated separately because the universal Mode 01 fuel-level PID (0x2F) is not consistently exposed or accurately scaled across manufacturers; the prototype therefore reads fuel level through a vehicle-specific adapter selected at install time, with the Toyota meter Mode 22 PID 21/29 channel used for the partner-fleet Toyota Hilux. The same telemetry cycle timestamps all acquired fields at the device and transmits them to the backend together with GPS, trip-state, and communication-channel metadata, so the anomaly detector interprets fuel changes in relation to vehicle operation rather than fuel percentage alone.

\subsection{Physical Interface}
The acquisition layer uses the vehicle OBD-II port and ISO 15765-4 transport over CAN. This enables the prototype to read fuel and vehicle context without tank-cap modification or sender-loom alteration. All OBD-II-compliant vehicles expose a standard set of Mode 01 PIDs over this transport, which the prototype uses for the universal contextual variables; manufacturer-specific PIDs (Toyota Mode 22, Hino J1939, Isuzu Mode 22) are layered on top of the same transport only for fuel-level acquisition.

The physical-interface choice is central to the study because it keeps the system non-invasive. Instead of inserting a new float sensor or cutting into the fuel sender wiring, the prototype reads the fuel-related signal already exposed by the vehicle electronics. This reduces installation risk, supports repeatability across fleet vehicles of the same model, and keeps the validation focused on whether the existing vehicle signal can be made accurate enough for monitoring.

\subsection{Vehicle-Specific Fuel-Level Adapter (Toyota Mode 22 PID 21/29)}
\label{sec:fuel_level_adapter}
Because the universal Mode 01 fuel-level PID (0x2F) was not reliably exposed by the partner-fleet ECU, the firmware falls back to a vehicle-specific fuel-level adapter selected at install time. For the Toyota Hilux deployed in this study the adapter is the Toyota meter Mode 22 command 0x21 0x29, which returns the raw fuel-sender voltage encoded at 0.5 L per ADC count without internal damping. The firmware polls this channel on a 1.5 s cadence and applies a retry policy when the adapter returns a negative acknowledgment or missing response. The adapter is implemented as a thin, replaceable layer over the universal OBD-II acquisition pipeline, so future deployments on other manufacturer ECUs (Hino J1939 SPN 96 for diesel heavy-duty units, Isuzu Mode 22, or any vehicle that fully implements Mode 01 PID 0x2F) require no changes to the filter cascade, telemetry schema, or anomaly-detection model.

The 1.5 s poll cadence was selected to capture fuel changes faster than the five-second telemetry publication interval while still leaving time for GPS parsing, display updates, and communication attempts. The retry policy prevents a single missing CAN response from becoming a false fuel event. Only values that pass validity checks are allowed to update the filtered fuel stream used by the HMI and backend.

\subsection{Co-acquired Mode 01 PIDs}

Table~\ref{tab:obd_pid_telemetry_mapping} lists the vehicle signals used by the prototype and shows where each signal is stored after ingestion. The column names are included only for traceability between the thesis results and the implemented database; the main purpose of the table is to show that fuel readings were interpreted together with speed, engine demand, motion context, and diagnostic status rather than as an isolated raw value.

\begin{table}[!htbp]
\centering
\caption{OBD-II PID to \texttt{telemetry\_logs} Column Mapping}
\label{tab:obd_pid_telemetry_mapping}
\scriptsize
\begin{tabular}{|p{0.22\textwidth}|p{0.33\textwidth}|p{0.35\textwidth}|}
\hline
\textbf{PID} & \textbf{Quantity} & \textbf{Telemetry column} \\
\hline
Mode 22 0x21 0x29 & Fuel sender voltage (0.5 L per ADC count) & \texttt{fuel\_level} (filtered) + \texttt{fuel\_source='mode22'} + \texttt{fuel\_confidence} \\
\hline
Mode 01 0x0C & Engine RPM & \texttt{engine\_rpm} \\
\hline
Mode 01 0x0D & Vehicle speed (km/h) & \texttt{speed} \\
\hline
Mode 01 0x04 & Calculated engine load (\%) & \texttt{engine\_load\_pct} \\
\hline
Mode 01 0x11 & Throttle position (\%) & \texttt{throttle\_pct} \\
\hline
Mode 01 0x10 & Mass air flow (g/s) & \texttt{maf\_gps} \\
\hline
Mode 01 0x05 & Engine coolant temperature (degC) & \texttt{coolant\_temp\_c} \\
\hline
Mode 01 0x1F & Time since engine start (s) & \texttt{engine\_runtime\_sec} \\
\hline
Mode 03 / 0x01 & Stored DTCs (count + presence) & \texttt{dtc\_present}, \texttt{dtc\_count} \\
\hline
\end{tabular}
\normalsize
\end{table}

Table~\ref{tab:obd_pid_telemetry_mapping} explains how raw vehicle signals become database-ready telemetry. Fuel level is the primary measurement for SO1 and SO4, while speed, RPM, load, throttle, mass airflow, coolant temperature, runtime, and diagnostic status provide context for deciding whether a fuel change is plausible. This is why the anomaly detector is described as context-aware: it does not evaluate fuel movement without also considering motion and engine state.

\subsection{Validity and Confidence Fields}
The firmware writes \texttt{fuel\_valid}, \texttt{fuel\_source}, and \texttt{fuel\_confidence} on every telemetry cycle. These fields allow the backend and dashboard to distinguish valid Mode 22 fuel readings from missing, stale, or fallback values.

These fields prevent the backend from treating all fuel values as equally reliable. For example, a current Mode~22 reading with high confidence can be used directly for alerts, while a stale or fallback value can be displayed with caution and excluded from model training or calibration windows. This protects the system from creating false anomalies during adapter timeouts or network replay.

\subsection{Timestamping}
Each telemetry record carries the original device-side timestamp as \texttt{original\_device\_\allowbreak timestamp}, while the backend records its own receipt time. Backend reconciliation uses server time to separate transmission delay from actual sampling time.

Dual timestamping is necessary because the communication subsystem deliberately permits delayed delivery. A row replayed from SPIFFS may arrive minutes after it was sampled, but its device timestamp still represents the true trip timeline. Chapter~6 therefore uses device time for sampling cadence and route reconstruction, and backend receipt time for communication latency and alert-arrival analysis.

\subsection{Field Reference Choice}
The manufacturer dashboard gauge is used as a supporting field reference for fuel-validation activity. This choice follows the delimitations in Chapter~\ref{ch:intro}, because direct dipstick validation and tank-cap opening were not permitted by the partner fleet operator.

The validation reference was therefore selected to match the allowed field procedure. Known-volume pump-gas events provide the volumetric reference, while the manufacturer dashboard confirms that the prototype is reading a fuel state consistent with the vehicle's own display. This is an honest operational validation rather than a laboratory calibration claim.

Because the validation reference is the legally calibrated trade-instrument fuel pump (Tables~\ref{tab:so1_pump_validation}--\ref{tab:so1_accuracy_summary}) rather than a Toyota-specific signal, the SO1 accuracy claim is portable to any future vehicle-specific adapter that produces a stable filtered fuel percentage; only the adapter in Section~\ref{sec:fuel_level_adapter} needs to be re-validated per ECU.

\section{Fuel Preprocessing Methodology}
\label{sec:fuel_preprocessing}

The raw OBD-II fuel reading was passed through a five-stage edge filter before publication. Stage~1 is a median window that rejects single-sample spikes caused by slosh and sender chatter. Stage~2 is a first-order infinite-impulse-response (IIR) exponential moving average (EMA) that smooths residual noise. Stage~3 is a state-dependent dead-band publisher that suppresses display jitter. Stage~4 is a direction-aware step rule that enforces the physical asymmetry of fuel: fuel can drop through combustion or removal, but can only rise through refueling. Stage~5 is a five-minute post-refuel quiescence window that prevents the published fuel percentage from tracking float-arm settling after a fill. The published output feeds both the HMI and the machine-learning pipeline; every subsequent feature in Section~\ref{sec:ml} derives from this value rather than the raw sender reading.

\begin{figure}[!htbp]
    \centering
    \includegraphics[width=0.95\textwidth]{Fuel-level_filter_cascade.png}
    \caption{Fuel-Level Filter Cascade}
    \label{fig:fuel_filter_cascade}
\end{figure}

\subsection{Sampling Rate and Notation}
\label{sec:fuel_notation}

Let $t \in \mathbb{Z}_{\geq 0}$ denote the discrete OBD sample index, with a fixed sampling interval $\Delta_s = 1.5$~s. The 1.5~s cadence is derived from three physical constraints. First, ISO~15765-4 imposes a minimum inter-request spacing of approximately 50~ms per PID, but the Toyota Mode~22 PID~21/29 fuel response on the partner-fleet Hilux was measured to take 60--120~ms per query, so any polling rate above approximately 5~Hz saturates the shared CAN bus and produces sporadic dashboard-warning flicker. Second, the physical maximum fuel-loss rate under idle combustion for a 2.4~L diesel engine is approximately $0.014$~L/min~$\approx 0.00029$~\%/s of the 80~L tank (SAE J1349 idle-combustion benchmark), so Nyquist coverage of legitimate dynamics requires only $\approx 0.001$~Hz; the 1.5~s cadence samples at approximately 30$\times$ the required rate, leaving substantial margin for burst-siphon detection. Third, the 1.5~s cadence times the 31-sample median window (Section~\ref{sec:fuel_median}) yields a 46.5~s memory that brackets the 5--20~s empirical slosh envelope measured during on-road cornering.

The following symbols are used throughout this section:
\begin{itemize}
    \item $x_t \in [0, F_{\max}]$: the raw Mode~22 sender reading at sample $t$, expressed as tank volume in litres.
    \item $m_t$: the output of Stage~1 (median filter) at sample $t$, in litres.
    \item $y_t$: the output of Stage~2 (EMA) at sample $t$, converted to percentage of tank fullness (\%).
    \item $p_t$: the \emph{published} fuel percentage at sample $t$~---~the value shown on the HMI and written to \texttt{telemetry\_logs.fuel\_level}.
    \item $p^{a}_{t-1}$: the last previously-published anchor value~---~the reference against which Stages~3 and~4 measure change.
    \item $s_t$: the vehicle speed at sample $t$ in km/h, from OBD-II Mode~01 PID~0x0D.
    \item $\mathbb{1}[\cdot]$: the indicator function, equal to 1 when the enclosed condition holds and 0 otherwise.
\end{itemize}

\subsection{Stage~1: Median Filter for Slosh Rejection}
\label{sec:fuel_median}

The Stage~1 output is the rank-order statistic
\begin{equation}
m_t = \operatorname{median}\bigl( x_{t-W+1},\, x_{t-W+2},\, \ldots,\, x_t \bigr), \qquad W = 31 \text{ samples} \approx 46.5\text{ s}
\label{eq:median_filter}
\end{equation}
where $W \in \mathbb{N}$ is the window length (odd, so the middle rank is unambiguous) and the operator returns the $((W+1)/2)$-th ordered value of the enclosed set. During warm-up, when the number of accumulated samples $n < W$, the HMI displays the placeholder \texttt{Loading} and no telemetry row is published; this prevents a single noisy first read from anchoring the display.

The choice $W = 31$ was reached empirically after the original $W = 15$ (approximately 22~s memory) failed on 2026-06-23 during on-road testing when a post-refuel float walked $45 \to 55 \to 54 \to 49$~\% over roughly 30~s. Characterisation of the 12 refuel events in the 27-trip corpus produced a 95th-percentile post-fill slosh duration of 42~s, implying $W \cdot \Delta_s \geq 42$~s and therefore $W \geq 28$; the next odd integer is 31. A larger window ($W = 61 \approx 90$~s) was also tested and rejected because it delayed real-burn response by 45~s or more during highway cruise, producing a display that drivers reported as \emph{frozen}. The rank-order statistic is retained instead of a mean because it has an optimal breakdown point of 50\%: a single 5~L slosh spike drags a mean but leaves the median unchanged as long as fewer than 15 of the 31 samples are contaminated (Rousseeuw \& Croux, 1993).

\subsection{Stage~2: First-Order IIR Exponential Moving Average}
\label{sec:fuel_ema}

The Stage~2 output is
\begin{equation}
y_t = \alpha \cdot m_t + (1 - \alpha) \cdot y_{t-1}, \qquad \alpha = 0.04, \quad y_0 = m_W
\label{eq:ema_filter}
\end{equation}
where $\alpha \in (0, 1]$ is the smoothing coefficient and $y_{t-1}$ is the internal filter state carried from the previous sample. Equation~\ref{eq:ema_filter} is a standard first-order infinite-impulse-response low-pass filter (Oppenheim \& Schafer, 2010). The equivalent time constant and the settling time to 95\% of a step input are
\begin{equation}
\tau = \frac{-\Delta_s}{\ln(1 - \alpha)} = \frac{-1.5}{\ln(0.96)} \approx 36.7 \text{ s}, \qquad t_{95} \approx 3\tau \approx 110 \text{ s}.
\label{eq:ema_time_constant}
\end{equation}

The choice $\alpha = 0.04$ was made through two independent lines of evidence. A power-spectral-density analysis of the post-median fuel trace from Trip~06 (a 55-minute highway cruise) showed a noise floor of $0.03$~\%$^2$ above $0.02$~Hz; setting the filter corner frequency $f_c = \alpha / (2\pi \Delta_s) = 0.0042$~Hz places that band at least 12~dB below the passband, which is the minimum attenuation needed for the Stage~3 dead-band to see a stable input. A parameter sweep over $\alpha \in \{0.01, 0.02, 0.04, 0.08, 0.15\}$ against the 27-trip corpus, minimising the mean absolute error between $y_t$ and a burn-rate ground truth constructed from odometer-normalised fuel-per-km, produced the minimum error at $\alpha = 0.04$; $\alpha = 0.08$ introduced approximately $3\times$ overshoot after refuels, and $\alpha = 0.02$ lagged real burn by more than four minutes. The 110~s settling time in Equation~\ref{eq:ema_time_constant} is deliberately absorbed by the warm-up rule: the 46.5~s median warm-up runs first, so the marginal user-visible wait added by the EMA is approximately 63.5~s. A driver expectation study using the System Usability Scale ($n = 8$) rated a \texttt{Loading} display of under two minutes as \emph{acceptable} (mean 4.1/5).

\subsection{Stage~3: State-Dependent Dead-Band Publisher}
\label{sec:fuel_deadband}

The Stage~3 output re-uses the previous published anchor unless the smoothed value has moved more than a motion-conditioned dead-band:
\begin{equation}
p_t =
\begin{cases}
y_t & \text{if } |y_t - p^{a}_{t-1}| \geq \theta_{\mathrm{db}}(s_t) \\[2pt]
p^{a}_{t-1} & \text{otherwise}
\end{cases}
\qquad
\theta_{\mathrm{db}}(s_t) =
\begin{cases}
0.5\,\% & \text{if } s_t > 0 \text{ within last 5 s (moving)} \\[2pt]
1.5\,\% & \text{otherwise (stationary)}
\end{cases}
\label{eq:deadband_filter}
\end{equation}
where $\theta_{\mathrm{db}}(\cdot)$ is the state-dependent dead-band threshold in percentage points and the anchor $p^{a}_{t-1}$ is the last value actually written to the display (not the last $y$).

The dual threshold reflects a physically-grounded noise asymmetry that was measured directly on the prototype. A six-hour overnight recording on TRCK-02 with the ignition off produced a standard deviation of sender drift of 0.42\%, with 99th-percentile excursion of 1.42\%. The 1.5\% stationary threshold puts drift below the dead-band with probability greater than 0.99. Moving noise was characterised on a 55-minute highway cruise after the EMA and produced a standard deviation of 0.11\% and a 99th-percentile of 0.38\%; the 0.5\% moving threshold covers 99\% of moving noise while still resolving the driver-visible burn rate on the Hilux of 0.17\%/min at 60~km/h~---~a value the 0.5\% dead-band allows to update the display every approximately three minutes, matching the human perception threshold for observing that fuel is falling.

\subsection{Stage~4: Direction-Aware Step Rule}
\label{sec:fuel_direction_aware}

The Stage~4 rule enforces the physical asymmetry that fuel cannot spontaneously rise while the truck is in motion. Let $\Delta_t = y_t - p^{a}_{t-1}$ denote the absolute percentage-point change between the current smoothed value and the anchor (this is a change, not a gradient and not $dF/dt$), and let $T_{\mathrm{park}}$ denote the cumulative time in seconds since the last observation of $s_t > 0$. The rule is:
\begin{equation}
p_t =
\begin{cases}
y_t & \text{if } \Delta_t \leq -2.5\,\% \quad\text{(down: real burn, always allowed)} \\[2pt]
y_t & \text{if } \Delta_t \geq +10\,\%\ \text{and}\ T_{\mathrm{park}} \geq 60\text{ s} \quad\text{(up: refuel event, gated)} \\[2pt]
p^{a}_{t-1} & \text{otherwise (blocked as slosh)}
\end{cases}
\label{eq:direction_aware}
\end{equation}

The $-2.5$\% down-step threshold was adapted from Barbado and Corcho (2022), whose Figure~4 places the 97.5th percentile of legitimate step-drops in a rule-based fuel-anomaly detector at approximately 2.3--2.6\%. This threshold was independently re-computed on the 27-trip corpus and produced a matching own-fleet 97.5th percentile of 2.31\%, so 2.5\% is slightly conservative: it rejects marginal outliers but never blocks a real burn event. The $+10$\% up-step threshold is derived from an audit of 340 SGSA fuel-log dispensations, in which the minimum recorded fill was 12~L on the 80~L tank (15\% of tank), so the 10\% floor provides a $1.5\times$ safety margin below the smallest legitimate fill. The 60~s stationary requirement corresponds to the 5th percentile of measured pump-transaction duration at four SGSA-approved fuel stations from the same log (median 42~s for driver bathroom breaks with the engine running, 5th-percentile of actual fills 62~s), which sets the floor just below the shortest legitimate fill while blocking driver-side false triggers.

\subsection{Stage~5: Post-Refuel Quiescence Window}

Once Stage~4 accepts an upward jump as a refuel event at sample $t^{r}$, the published value $p_t$ is held at $y_{t^r}$ for the next \texttt{REFUEL\_QUIESCENCE\_MS}~$= 300{,}000$~ms (five minutes). This suppresses the tail of float-arm settling that normally follows a fill event: the sender arm rebounds through a damped oscillation with an observed 90\% settling time of 180--240~s on the partner-fleet Hilux, and the 300~s window covers this transient conservatively. During the quiescence window, both upward and downward changes are ignored because true consumption over five minutes at highway cruise (approximately 0.85\% on 80~L) is dwarfed by the settling amplitude of $\pm$3--5\%.

\subsection{Effect on the Machine-Learning Pipeline}

The practical effect of Equations~\ref{eq:median_filter}--\ref{eq:direction_aware} together with Stage~5 is that the published percentage $p_t$ is deterministic, non-negative in its change during motion, and monotone downward outside of confirmed refuel windows. This matters for both SO1 and SO4. SO1 evaluates whether the displayed level remains accurate against pump-validation events; the cascade removes the sender-jitter that would otherwise obscure this comparison. SO4 depends on the cascade removing one-sample spikes so that the Isolation Forest sees fuel changes that reflect real state transitions rather than sender chatter. Every feature engineered in Section~\ref{sec:ml} (in particular \texttt{fuel\_delta}, \texttt{fuel\_per\_km}, \texttt{fuel\_rate\_per\_hour}, and the \texttt{fuel\_drop\_*} family) derives from $p_t$, so the cascade is the foundation on which the model's precision claim rests.

\section{Communication and Offline Sync Methodology}
\label{sec:communication}

\subsection{Priority Chain Rationale}
LTE was treated as the primary wide-area data path because it offers the lowest steady-state latency and highest payload throughput among the available wide-area options. LoRa was retained as the depot/yard and long-range telemetry path because of its multi-kilometre line-of-sight reach at very low power. GSM 2G was retained as the deeper-fallback for 2G-only coverage pockets on the Batangas and Subic routes. WiFi was retained as a depot or driver-tether path. SPIFFS offline-buffer replay is the final fallback that guarantees no telemetry record is silently lost when all four real-time channels are unavailable.

\subsection{Formal Channel-Selection Model}
\label{sec:comm_channel_selection}

Let $\mathcal{C} = \{c_1, c_2, c_3, c_4, c_5\}$ denote the ordered set of transport channels indexed by priority, with $c_1 = \text{LoRa}$, $c_2 = \text{LTE}$, $c_3 = \text{GSM~2G}$, $c_4 = \text{WiFi}$, and $c_5 = \text{SPIFFS}$. Let $\mathrm{Online}(c, t) \in \{0, 1\}$ be the channel-specific liveness predicate defined below. Each telemetry cycle $t$ selects the highest-priority live channel:
\begin{equation}
r(t) = \operatorname*{argmin}_{c_i \in \mathcal{C}} \bigl\{ i \;\bigm|\; \mathrm{Online}(c_i, t) = 1 \bigr\}
\label{eq:route_selection}
\end{equation}
where $r(t)$ is the channel index selected for cycle $t$. Because $c_5$ (SPIFFS) is always writable, the argmin is always defined. Equation~\ref{eq:route_selection} is re-evaluated every $\Delta_T = 5{,}000$~ms (\texttt{TELEMETRY\_INTERVAL\_MS}); \emph{no channel is locked in for consecutive cycles.} If LoRa is selected at cycle $t$ and fails during $t+1$, LTE is tried at $t+1$ without waiting for a lock-out interval. This distinguishes the prototype from a session-oriented cellular router and is what preserves the total system delivery rate reported in Chapter~\ref{ch:result_discuss}.

\subsection{Per-Channel Connection Detection}
\label{sec:comm_online_predicates}

Each channel implements $\mathrm{Online}(c, t)$ differently because the physical layer defines what \emph{connected} means for each medium. The four active predicates are:

\begin{itemize}
    \item \textbf{WiFi} (four-way handshake + Layer-7 probe):
        \begin{equation}
        \mathrm{Online}(\text{WiFi}, t) = \mathbb{1}\bigl[\, \text{WiFi.status()} = \texttt{WL\_CONNECTED} \,\bigr] \cdot \mathbb{1}\bigl[\, \text{HTTP GET}\ /\text{health}\ \text{returned 2xx within 30~s} \,\bigr]
        \label{eq:online_wifi}
        \end{equation}
        The Layer-7 probe is required because captive portals and DNS failures produce a \texttt{WL\_CONNECTED} state without backend reachability. The 30~s probe interval was selected because the ESP32 WiFi module draws approximately 90~mA active versus 20~mA idle, and probing every 30~s adds approximately 0.2\% duty $\approx 0.02$~W average~---~negligible against the DHCP lease-refresh cadence.

    \item \textbf{LTE} (PDP context activation + signal-strength refresh):
        \begin{equation}
        \mathrm{Online}(\text{LTE}, t) = \mathbb{1}\bigl[\, \texttt{gsmReady} = \text{true after AT+CGACT=1,1 OK} \,\bigr] \cdot \mathbb{1}\bigl[\, \text{CSQ poll returned }+\text{CSQ:}\text{ within last }3\text{ s} \,\bigr]
        \label{eq:online_lte}
        \end{equation}
        The CSQ (received signal-strength) command is issued every $T_{\text{CSQ}} = 1{,}500$~ms; two consecutive misses (3~s) declare the modem offline. The 1.5~s cadence aligns with the OBD poll so both share the same Core~0 scheduler tick.

    \item \textbf{GSM~2G:} the A7670E modem falls back to EDGE internally when LTE registration fails. The predicate uses the same CSQ liveness envelope as LTE, so the same Equation~\ref{eq:online_lte} applies with the modem-reported RAT tag distinguishing the two rows in the diagnostic log.

    \item \textbf{LoRa} (per-packet acknowledgment with backoff):
        \begin{equation}
        \mathrm{Online}(\text{LoRa}, t) = \mathbb{1}\bigl[\, (t - t^{\mathrm{ack}}) < 90{,}000\text{ ms} \,\bigr] \cdot \mathbb{1}\bigl[\, (t - t^{\mathrm{fail}}) > 60{,}000\text{ ms} \,\bigr]
        \label{eq:online_lora}
        \end{equation}
        where $t^{\mathrm{ack}}$ is the timestamp of the most recent gateway acknowledgment and $t^{\mathrm{fail}}$ is the timestamp of the most recent 2~s ACK-timeout event. A single packet is transmitted with an ACK-timeout of 2~s. The 2~s value bounds the measured 99th-percentile round-trip of 1.31~s at spreading factor~7, which contains a 62~ms airtime for a 51-byte payload times two (round-trip) plus a worst-case measured 850~ms gateway processing delay. The 60~s failure back-off prevents thrashing when the truck rolls out of gateway range (60~s $\approx$ typical urban block length at 40~km/h). The 90~s liveness window is deliberately larger than the back-off so that a single missed ACK never disqualifies LoRa while the truck is momentarily blocked by traffic.
\end{itemize}

\subsection{Background Polling Cost and Cached-Handshake Argument}
\label{sec:comm_polling_cost}

A panel question raised during preliminary review was whether the multi-channel handshake logic introduces per-packet latency. It does not, because the readiness predicates in Equations~\ref{eq:online_wifi}--\ref{eq:online_lora} are cached flags refreshed by inexpensive background polls rather than re-executed handshakes before every send. The cumulative CPU cost is
\begin{equation}
C_{\mathrm{poll}} = \frac{t^{\mathrm{CSQ}}_{\mathrm{cmd}}}{T_{\mathrm{CSQ}}} + \frac{t^{/\mathrm{health}}_{\mathrm{cmd}}}{T_{/\mathrm{health}}} + 0
\label{eq:polling_cost}
\end{equation}
where $t^{\mathrm{CSQ}}_{\mathrm{cmd}} \approx 10$~ms, $T_{\mathrm{CSQ}} = 1{,}500$~ms, $t^{/\mathrm{health}}_{\mathrm{cmd}} \approx 200$~ms, $T_{/\mathrm{health}} = 30{,}000$~ms, and the LoRa contribution is zero because LoRa liveness is refreshed passively by the ACK of the previous packet. Substituting yields $C_{\mathrm{poll}} \approx 6.67 + 6.67 = 13.3$~ms/s, or approximately 1.3\% of the Core~0 duty cycle.

\subsection{Per-Packet Latency Decomposition}
\label{sec:comm_latency_decomp}

The end-to-end telemetry latency for a warm session is decomposed as
\begin{equation}
T_{\mathrm{total}}(t) = T_{\mathrm{handshake}}^{\mathrm{cached}} + T_{\mathrm{TX}}(c, |b|) + T_{\mathrm{RTT}}^{\mathrm{backend}}
\label{eq:latency_decomp}
\end{equation}
where $T_{\mathrm{handshake}}^{\mathrm{cached}} \approx 0$ for a warm session (Section~\ref{sec:comm_polling_cost} above), $T_{\mathrm{TX}}(c, |b|)$ is the airtime of a payload of size $|b|$ on channel $c$, and $T_{\mathrm{RTT}}^{\mathrm{backend}}$ is the Supabase remote-procedure-call round-trip time. For the LoRa 51-byte payload at SF7, $T_{\mathrm{TX}} = 62$~ms; for WiFi the airtime is negligible against the WAN round trip. The 7.2~s WiFi median and 7.4~s LoRa median latencies reported in Chapter~\ref{ch:result_discuss} are therefore dominated by $T_{\mathrm{RTT}}^{\mathrm{backend}}$~---~not by the channel-selection logic. A counterfactual controlled measurement confirmed this: forcing a full TLS handshake before every send added 300--500~ms of measured overhead, which motivates the cached-flag design.

\subsection{Communication Module Comparison}

\begin{table}[!htbp]
\centering
\caption{Communication Module Comparison}
\label{tab:communication_module_comparison}
\scriptsize
\setlength{\tabcolsep}{2pt}
\renewcommand{\arraystretch}{1.15}
\begin{tabular}{|p{0.08\textwidth}|p{0.12\textwidth}|p{0.13\textwidth}|p{0.10\textwidth}|p{0.25\textwidth}|p{0.07\textwidth}|p{0.13\textwidth}|}
\hline
\textbf{Module} & \textbf{Technology} & \textbf{Hardware} & \textbf{Bandwidth} & \textbf{Range} & \textbf{Power} & \textbf{Role} \\
\hline
WiFi & 802.11 b/g/n & ESP32 built-in & approx. 10 Mbps & approx. 50 m & Low & Depot and driver hotspot \\
\hline
LTE & LTE Cat-1 & A7670E & approx. 10 Mbps & Cell coverage & Medium & Primary cellular path \\
\hline
GSM 2G & EDGE & A7670E fallback & approx. 80 kbps & Wider than LTE & Medium & Cellular fallback \\
\hline
LoRa & SX1278, 433 MHz, SF\glsadd{SF}7, BW\glsadd{BW} 125 kHz & RA-02, 17 dBm TX & approx. 5.5 kbps & approx. 50 m urban NLOS\glsadd{NLOS} at 97\% measured; approx. 1 km open LOS\glsadd{LOS} at 60\% measured; up to 10 km LOS per SX1278 datasheet & Very low & Primary telemetry path (intelligent traffic-steering chain) \\
\hline
Offline & - & SPIFFS flash & - & - & None & Outage replay \\
\hline
\end{tabular}
\normalsize
\end{table}

The LoRa range column is reported as two values: the operational urban NLOS reach measured in this study and the line-of-sight reach published in the SX1278 datasheet. LoRa is the highest-priority transport in the intelligent-traffic-steering chain (LoRa $\rightarrow$ LTE $\rightarrow$ GSM 2G $\rightarrow$ WiFi $\rightarrow$ SPIFFS buffer); the cellular path provides coverage of the urban-NLOS regime past approximately 50 m where LoRa delivery degrades.

Table~\ref{tab:communication_module_comparison} is not a generic radio comparison; it explains why every communication option remains in the prototype. LoRa gives low-power depot and yard coverage, LTE provides the main wide-area path, GSM~2G provides a legacy fallback in fringe coverage, WiFi supports depot or driver-hotspot use, and SPIFFS protects the data when no radio path is available. The later SO2 tests measure the actual performance of this chain.

\subsection{Communication Method Comparison}

\begin{figure}[!htbp]
    \centering
    \includegraphics[width=0.8\textwidth]{Communication_Method_Comparison_radar.png}
    \caption{Communication Method Comparison}
    \label{fig:communication_method_radar}
\end{figure}

Figure~\ref{fig:communication_method_radar} compares LoRa, LTE, GSM 2G, WiFi, and SPIFFS on range, bandwidth, power efficiency, Philippine coverage, infrastructure independence, and data integrity using a five-point normalized scale.

The radar chart summarizes the trade-off that motivates the fallback design. No single channel dominates all criteria: LoRa is efficient and infrastructure-light but low bandwidth, LTE has broad operational usefulness but depends on carrier coverage, GSM~2G is slow but useful in legacy pockets, WiFi is fast but short-range, and SPIFFS has perfect local retention but no real-time delivery. The system therefore treats communication as a prioritized set of options rather than as one fixed modem path.

\subsection{Telemetry Payload Summary}

Each telemetry payload carries the identifiers needed for synchronization, the filtered fuel reading and confidence flags, GPS position and source metadata, OBD-II context variables, communication-channel tags, trip-state counters, and distance deltas. This compact payload design lets the backend deduplicate records, separate live from replayed telemetry, and run anomaly detection without requiring the dashboard to parse device-specific state. The full JSON payload schema is provided in Appendix~\ref{ch:appendix_technical_details}.

\subsection{Offline-Buffer Fall-Through Flowchart}

Figure~\ref{fig:offline_buffer_fallthrough} shows the decision tree executed by the Core~0 worker each telemetry cycle, tracing how a record is routed through the priority chain and appended to the SPIFFS buffer when all online channels are unavailable.

\begin{figure}[!htbp]
    \centering
    \includegraphics[width=0.9\textwidth]{Offline-buffer_fall-through_flowchart.png}
    \caption{Offline-Buffer Fall-Through Flowchart}
    \label{fig:offline_buffer_fallthrough}
\end{figure}

\noindent
The flowchart is read from top to bottom as a reliability mechanism. A telemetry row is first offered to the preferred live channel; if that channel fails, the same row is attempted on the next path without changing its sequence identifier. Only after all live attempts fail is the record stored locally. This design supports later deduplication because replayed records still carry their original event identity.

\subsection{Fall-through Policy}
Each telemetry record is attempted on the highest-priority available channel: LoRa first if a gateway uplink is in range, otherwise LTE, then GSM 2G, then WiFi. If an attempt fails because of a non-2xx response from the backend or a modem error, the firmware falls through to the next channel without duplicating the \texttt{seq} or \texttt{event\_id}. If all real-time channels fail, the record is appended to the SPIFFS ring buffer and replayed once any channel recovers. The replayed record is tagged \texttt{channel\_\allowbreak used = 'offline\_\allowbreak replay'} so the backend can distinguish live from replayed data.

\subsection{Security}
Each telemetry payload is signed with an HMAC\glsadd{HMAC} digest before transmission. All cellular and WiFi data paths use HTTPS\glsadd{HTTPS} with TLS, protecting telemetry records against tampering and eavesdropping in transit. LoRa payloads are signed at the application layer, as the LoRa physical layer does not provide built-in encryption. These measures ensure that data integrity and confidentiality are maintained across all communication channels in the priority chain.

\subsection{LoRa Range Test Protocol}
The LoRa physical-layer range test was conducted with the SX1278 module configured at seventeen-dBm transmit power, four-hundred-thirty-three-megahertz centre frequency, spreading factor seven, and one-hundred-twenty-five-kilohertz bandwidth. One hundred telemetry-sized packets were transmitted from the in-vehicle node to a fixed gateway at each of the open line-of-sight distances of ten, fifty, one hundred, two hundred and fifty, five hundred, one thousand, two thousand, and five thousand metres, and at the urban non-line-of-sight distances of one hundred, two hundred and fifty, and five hundred metres. The received-signal-strength indicator, the signal-to-noise ratio, and the successfully-received-packet count were recorded at the gateway for each distance band.

\subsection{GSM / LTE Drive-test Protocol}
Carrier-reported signal quality was logged via the AT+CSQ\glsadd{CSQ} command issued to the A7670E modem at five locations of operational interest: the depot, a Metro Manila urban environment, a highway corridor, a tunnel or underpass selected to trigger the buffer-activation path, and a provincial route in Batangas exhibiting two-gigahertz-only coverage. The per-packet upload latency was computed at the backend as the difference between the device-side \texttt{sent\_at} timestamp and the backend-side receipt timestamp; the fiftieth- and ninety-fifth-percentile latencies are reported per location alongside the successful-upload fraction.

\subsection{Packet-delivery and Buffering Protocol}
A drive route was executed that traversed three distinct connectivity regions: an urban region with strong LTE coverage, a highway region with marginal LTE coverage, and a confirmed no-signal dead zone. One hundred packets were generated at each stage; the per-stage outcome (sent on first attempt, sent after a retry, buffered to the on-device SPIFFS log, replayed once connectivity returned, or lost) was tabulated from the firmware serial-output log and the backend-side telemetry receipt log.

Together, the LoRa range test, GSM/LTE drive test, and packet-delivery route test separate three different communication questions: how far the radio link can reach, how cellular coverage behaves in real service areas, and whether the overall system preserves data when coverage disappears. This separation makes the SO2 results interpretable because a latency problem, a radio-range problem, and a buffering problem are not treated as the same failure mode.

\section{GPS Acquisition and Fusion Methodology}
\label{sec:gps_pipeline}

The GPS pipeline provides three independent products: instantaneous position for driver navigation, an integrated trip-distance estimate for fuel-economy features, and an independent speed source that cross-checks the OBD-II speed PID. All three products feed both the human-machine interface and the machine-learning pipeline, and each has its own timing and validity envelope.

\subsection{Physical Interface and NMEA Extraction}
The receiver is a u-blox NEO-6M GPS module (50 channels, $-161$~dBm tracking sensitivity, datasheet cold-start of at most 27~s) connected to the ESP32 by UART at 9600~baud on GPIO16 (RX) and GPIO17 (TX). The NMEA byte stream is drained on every scheduler tick and fed to the TinyGPSPlus embedded parser, which maintains a finite-state machine over four sentence types: \$GPRMC/\$GNRMC (position, ground speed, heading, UTC), \$GPGGA/\$GNGGA (fix quality, satellites used, HDOP, altitude), \$GPGSA/\$GNGSA (fix type 1/2/3D), and \$GPGSV/\$GNGSV (satellites in view, for diagnostics only). Each successful fix produces the tuple
\begin{equation}
f_t = \bigl(\varphi_t,\; \lambda_t,\; v^{\mathrm{gps}}_t,\; h_t,\; \theta^{\mathrm{c}}_t,\; \mathrm{HDOP}_t,\; n^{\mathrm{sats}}_t,\; \tau_t \bigr)
\label{eq:gps_tuple}
\end{equation}
where $\varphi_t$ and $\lambda_t$ are the WGS-84 latitude and longitude in degrees, $v^{\mathrm{gps}}_t$ is the ground speed in km/h obtained from the RMC sentence (independent of the OBD-II speed PID), $h_t$ is the ellipsoidal altitude in metres, $\theta^{\mathrm{c}}_t$ is the ground course in degrees true (used for the driver-facing navigation arrow but not passed to the machine-learning pipeline), $\mathrm{HDOP}_t$ is the horizontal dilution of precision (dimensionless), $n^{\mathrm{sats}}_t$ is the count of satellites contributing to the fix, and $\tau_t$ is the UTC epoch timestamp.

\subsection{Haversine Great-Circle Distance}
\label{sec:gps_haversine}

The distance between two consecutive valid fixes is computed with the haversine great-circle formula
\begin{equation}
a = \sin^{2}\!\left(\frac{\Delta\varphi}{2}\right) + \cos(\varphi_1)\cdot\cos(\varphi_2)\cdot\sin^{2}\!\left(\frac{\Delta\lambda}{2}\right), \qquad d_{\mathrm{seg}} = 2R\cdot\arcsin\!\bigl(\sqrt{a}\bigr)
\label{eq:haversine}
\end{equation}
where $\varphi_1, \varphi_2$ and $\lambda_1, \lambda_2$ are the latitude and longitude of the two consecutive fixes in radians, $\Delta\varphi = \varphi_2 - \varphi_1$, $\Delta\lambda = \lambda_2 - \lambda_1$, $R = 6{,}371$~km is the mean Earth radius under the WGS-84 spherical approximation, and $d_{\mathrm{seg}}$ is the great-circle chord length in kilometres. The intermediate quantity $a$ is dimensionless.

\subsection{Speed-Gated Integration}
\label{sec:gps_integration}

The trip-level distance accumulator applies two gating conditions before adding a segment:
\begin{equation}
D^{\mathrm{gps}}_t = D^{\mathrm{gps}}_{t-1} + d_{\mathrm{seg}} \cdot \mathbb{1}\bigl[\,v^{\mathrm{gps}}_t > 3\text{ km/h}\,\bigr] \cdot \mathbb{1}\bigl[\,d_{\mathrm{seg}} < 0.5\text{ km}\,\bigr]
\label{eq:gps_integration}
\end{equation}
where $D^{\mathrm{gps}}_t$ is the cumulative GPS-integrated distance in kilometres at sample $t$. Both indicators must evaluate to 1 for the segment to be accumulated.

The 3~km/h speed floor is drawn from the NEO-6M datasheet: the receiver's stationary CEP-50 velocity noise is $\pm$1.5~km/h, so 3~km/h~$= 2\times$ the noise floor. This rejects the approximately 0.0003~km per sample multipath drift that otherwise adds approximately 1~km per hour of parked-truck distance. The 0.5~km segment cap is derived from a legitimate-motion upper bound: at the Philippine expressway speed cap of 100~km/h $= 27.8$~m/s, one second at 1~Hz fix cadence produces at most 28~m~$= 0.028$~km, so 0.5~km provides an approximately $18\times$ headroom above the fastest legitimate segment and rejects only physically impossible teleport jumps (typically a lost-then-reacquired fix that appears approximately 1~km away after a tunnel).

\subsection{OBD-II / GPS Distance Fusion}
\label{sec:gps_fusion}

Distance is fused across three tiers of decreasing precision. The backend selects the fused distance per telemetry cycle as
\begin{equation}
\Delta d^{\mathrm{OBD}}_t =
\begin{cases}
o_t - o_{t-1} & \text{if Toyota Mode~22 PID DE0F valid} \\[2pt]
v^{\mathrm{obd}}_t \cdot \Delta_T & \text{otherwise (OBD speed integration)}
\end{cases}
\label{eq:obd_delta}
\end{equation}
\begin{equation}
\Delta d^{\mathrm{fused}}_t =
\begin{cases}
\Delta d^{\mathrm{GPS}}_t & \text{if } \Delta d^{\mathrm{GPS}}_t > 0 \\[2pt]
\Delta d^{\mathrm{OBD}}_t & \text{otherwise}
\end{cases}
\label{eq:fusion_rule}
\end{equation}
where $o_t$ is the cumulative OBD-II odometer reading in kilometres, $v^{\mathrm{obd}}_t$ is the OBD-II vehicle speed in km/h, $\Delta_T = 5$~s~$= 5/3600$~h is the telemetry cycle length, and $\Delta d^{\mathrm{GPS}}_t$ is the change in $D^{\mathrm{gps}}$ over the last telemetry cycle. The fused value $\Delta d^{\mathrm{fused}}_t$ is written to \texttt{telemetry\_logs.distance\_fused\_delta\_km} and is what the downstream features \texttt{fuel\_per\_km} and \texttt{distance\_per\_fuel} consume.

GPS is preferred whenever it reports a valid non-zero segment because OBD-II speed integration drifts by approximately 3\% over a 30-minute trip: OBD PID~0x0D returns quantised 1~km/h values, so integration over even a modest number of samples accumulates a bounded bias. The GPS haversine sum in Equation~\ref{eq:haversine} is memoryless~---~each segment stands on its own, so no such drift accumulates. The 3~km/h speed gate from Equation~\ref{eq:gps_integration} keeps GPS from over-reporting during idle windows.

\subsection{Freshness and Fix-Loss Envelope}
\label{sec:gps_freshness}

Two independent staleness rules operate in parallel. On the backend,
\begin{equation}
\mathrm{Fresh}(f_t) = \mathbb{1}\bigl[\,(t_{\mathrm{now}} - \tau_t) \leq T_{\mathrm{fresh}}\,\bigr], \qquad T_{\mathrm{fresh}} = 15{,}000\text{ ms}
\label{eq:gps_fresh}
\end{equation}
where $t_{\mathrm{now}}$ is the backend wall-clock time at insert time and $\tau_t$ is the GPS UTC timestamp from Equation~\ref{eq:gps_tuple}. The value $T_{\mathrm{fresh}} = 3 \cdot \Delta_T$ tolerates one missed telemetry cycle before the fix is considered stale. A separate window $T_{\mathrm{merge}} = 5{,}000$~ms controls when a mobile-phone GPS report and an embedded GPS report are merged onto the same telemetry row; $T_{\mathrm{merge}} = \Delta_T$ so contemporaneous mobile and embedded GPS from the same telemetry cycle land on the same row.

On the edge, a fix is considered lost after
\begin{equation}
\text{gpsFix} = \text{false} \quad \text{when} \quad (\text{millis()} - t_{\mathrm{lastFix}}) > \texttt{GPS\_FIX\_TIMEOUT\_MS}, \quad \texttt{GPS\_FIX\_TIMEOUT\_MS} = 5{,}000
\label{eq:gps_fix_timeout}
\end{equation}
at which point the HMI switches to the \texttt{GPS: LOST} indicator. The 5~s timeout equals one telemetry cycle, so the driver-visible loss indication and the backend's staleness decision are synchronised.

\section{Backend and Database Methodology}
\label{sec:backend}

\subsection{Backend Stack}
The backend stack consists of Node/Express hosted on Render and Supabase Postgres as the managed relational database. This stack receives device telemetry, manages trip state, runs alert rules, invokes the anomaly-detection service, and exposes synchronized state to the HMI, mobile driver application, and administrative dashboard.

\subsection{Database Schema Summary}

\begin{table}[!htbp]
\centering
\caption{Supabase Postgres Schema}
\label{tab:supabase_postgres_schema}
\scriptsize
\begin{tabular}{|p{0.20\textwidth}|p{0.36\textwidth}|p{0.34\textwidth}|}
\hline
\textbf{Record group} & \textbf{Purpose} & \textbf{Key fields} \\
\hline
Users & Accounts for administrators, fleet managers, managers, and drivers & User identifier, role, assigned fleet-vehicle identifier \\
\hline
Fleet vehicles & Vehicle profile and maintenance state & Vehicle identifier, fleet-vehicle code, plate number, current odometer, status \\
\hline
Trip sessions & Active / paused / ended trips with destination, route polyline, navigation state, and mobile-GPS fields & Trip identifier, fleet-vehicle identifier, driver identifier, trip status, distance, assigned destination, route steps, mobile location \\
\hline
Telemetry logs & Per-cycle fuel, GPS, OBD context, and anomaly flag & Fuel level, location, speed, odometer, communication channel, anomaly flag, anomaly score, model source \\
\hline
Archived telemetry logs & Older trip telemetry beyond retention window & Telemetry fields plus archive timestamp and archive reason \\
\hline
Trip summaries & Per-trip aggregates refreshed by the backend & Total distance, total alerts, total anomalies, average fuel level \\
\hline
Alerts & Rest, maintenance, overspeed, low-fuel, and fuel-anomaly events & Alert type, severity, resolution status, trip identifier \\
\hline
Latency logs & SO2 latency tracking between device send time and backend receipt time & Latency, communication channel, retry state, send time, receipt time \\
\hline
\end{tabular}
\normalsize
\end{table}

Table~\ref{tab:supabase_postgres_schema} summarizes the live database entities used by the system. The schema is organized around trip sessions because trips are the unit that bind a driver, a fleet vehicle, fuel readings, GPS rows, alerts, latency records, and final summaries. Keeping these entities relational rather than storing them as unstructured blobs allows the dashboard to filter by trip, vehicle, alert type, driver, and communication channel while preserving traceability for Chapter~6 evaluation.

\subsection{Entity-Relationship Diagram}

\begin{figure}[!htbp]
    \centering
    \includegraphics[width=0.95\textwidth]{fig_5_6_erd.png}
    \caption{ERD Diagram}
    \label{fig:erd_diagram}
\end{figure}

Figure~\ref{fig:erd_diagram} presents the Entity-Relationship Diagram (ERD\glsadd{ERD}) of the system's Supabase PostgreSQL database, containing eight interrelated record groups for fleet vehicles, users, trip sessions, telemetry logs, archived telemetry logs, trip summaries, alerts, and latency logs. The trip-session record serves as the main hub for the relationships among these groups, linking fleet vehicles and drivers. Other records reference trips to relate telemetry, alerts, summaries, and latency measurements to a specific trip session. This data structure is designed to maintain referential integrity throughout the schema and support scalable access to the HMI dashboard, mobile app, and backend analytic pipeline.

\subsection{Cross-Interface State Synchronization}

\begin{figure}[!htbp]
    \centering
    \includegraphics[width=0.95\textwidth]{fig_5_6b_state_sync.png}
    \caption{Cross-Interface State Synchronization}
    \label{fig:cross_interface_state_sync}
\end{figure}

Figure~\ref{fig:cross_interface_state_sync} illustrates the cross-interface synchronization flow of the system, depicting how data propagates across the HMI, backend, database, and dashboard layers in response to driver actions. When the driver initiates a trip via the HMI or mobile app, the backend creates an active trip record in \texttt{trip\_sessions}, which simultaneously triggers the fleet-vehicle card display to activate on the dashboard. Telemetry data is transmitted every 5000 ms as a cycle during a trip. Each cycle is processed and anomaly-detected by the backend before it is recorded in \texttt{telemetry\_logs} and displayed as alerts in the dashboard. Trip status changes are also reflected throughout the system through \texttt{trip\_sessions.\allowbreak trip\_status} for Rest, Resume, or End trip states. This mechanism ensures that all interfaces, including the mobile PWA and HMI poll-back, remain synchronized with the current session state. Throughout this process, \texttt{trip\_sessions} functions as the canonical source of trip data and the authoritative state controller for the entire architecture.

\subsection{Active-trip Uniqueness Enforcement}
Duplicate active-trip sessions across the three interfaces are prevented by partial unique indexes on the \texttt{trip\_sessions} table. Two indexes enforce the constraint: one restricts each fleet vehicle to at most one active trip at any time, and a second restricts each driver to at most one active trip at any time. Both indexes are conditioned on \texttt{trip\_status = 'active'} so that historical records with ended or paused status are not affected. Any attempt to open a second concurrent active trip for the same fleet vehicle or driver is rejected at the database layer before reaching application logic, eliminating race conditions that could arise when a driver starts a trip from both the HMI and the mobile web application simultaneously.

\subsection{Odometer and Trip-distance Methodology}
The fleet vehicles used in the prototype display lifetime mileage on the physical dashboard, but that value is not exposed through the OBD-II adapter available to the system. Although the firmware probes the manufacturer-specific odometer data identifiers, the fleet does not return a usable lifetime odometer value. As a result, the per-record odometer field is either empty or a trip-scoped, speed-integrated estimate generated by the embedded HMI, and it is not treated as an absolute lifetime odometer reading.

For this reason, backend lifetime mileage is computed from completed trip distance rather than from the raw OBD-II odometer field. Each active trip continually updates its accumulated trip distance using a distance-estimation routine that takes the greater of a capped speed-by-time integration and a filtered GPS haversine distance, rejecting stationary GPS jitter below a movement threshold. When a trip transitions to the ended state, the backend adds the finalised trip distance to the fleet vehicle's cumulative odometer exactly once and records the time at which the accumulation occurred. This design ensures that cumulative mileage remains consistent and is never double-counted, even when telemetry records arrive out of order or are replayed from the offline buffer.

\subsection{Alert Rule Implementation}
\label{sec:alerts}

\begin{table}[!htbp]
\centering
\caption{Alert Rule Thresholds Implemented}
\label{tab:alert_rule_thresholds}
\scriptsize
\setlength{\tabcolsep}{2pt}
\renewcommand{\arraystretch}{1.15}
\begin{tabular}{|p{0.15\textwidth}|p{0.27\textwidth}|p{0.20\textwidth}|p{0.28\textwidth}|}
\hline
\textbf{Alert} & \textbf{Threshold} & \textbf{Trigger source} & \textbf{Hardware action} \\
\hline
Rest reminder & distance $\geq$ 300 km AND drive time $\geq$ 4 h & backend \texttt{/device/alerts} & HMI buzzer, screen alert, \texttt{Take Rest} button \\
\hline
Maintenance (oil) & trip distance $\geq$ 5,000 km & backend trip aggregation & Dashboard banner + planned email/push \\
\hline
Overspeed & speed $\geq$ 100 km/h & telemetry stream & HMI flash and log \\
\hline
Fuel anomaly & ML flag from Isolation Forest or Matrix Profile & ML anomaly service & Dashboard alert card \\
\hline
\end{tabular}
\normalsize
\end{table}

Table~\ref{tab:alert_rule_thresholds} defines the operational thresholds used by the backend rule engine. These rules are deliberately simple and auditable: rest, maintenance, overspeed, and low-fuel alerts are deterministic, while fuel-anomaly alerts originate from the ML service and pass through the contextual gate. This distinction matters because deterministic safety rules can be validated by threshold crossing, whereas fuel-anomaly rules require the model-performance analysis in SO4.

\begin{table}[!htbp]
\centering
\caption{Extended Safety Implementation Status}
\label{tab:extended_safety_status}
\scriptsize
\begin{tabular}{|p{0.23\textwidth}|p{0.16\textwidth}|p{0.51\textwidth}|}
\hline
\textbf{Safety alert} & \textbf{Status} & \textbf{Implementation note} \\
\hline
Overspeed & Implemented & Backend rule at 100 km/h threshold against field-test telemetry speed \\
\hline
Rest reminder & Implemented & SO5-b spec is 300 km / 6 h; implemented at 4 h per LTO driver-fatigue reading \\
\hline
Maintenance (5,000 km) & Implemented & Backend trip-aggregation rule against cumulative odometer \\
\hline
Harsh acceleration / braking & Partial & Acceleration rate computed during preprocessing; no alert rule yet \\
\hline
Lane departure & Out of scope & Requires camera and IMU - future work \\
\hline
Fatigue / drowsiness & Out of scope & Requires camera with face detection - future work \\
\hline
\end{tabular}
\normalsize
\end{table}

Table~\ref{tab:extended_safety_status} separates implemented alerts from partial or out-of-scope safety concepts. Harsh acceleration/braking is marked partial because the feature pipeline computes the signal but no user-facing rule has been deployed. Lane departure and fatigue/drowsiness are excluded because they require sensing hardware not present in the prototype. This prevents the thesis from implying safety functions that the installed hardware cannot support.

\subsection{Archival Pipeline}
Telemetry records older than the active-retention window are moved from \texttt{telemetry\_logs} to \texttt{archived\_\allowbreak telemetry\_logs} by a backend archival pipeline. The archived table retains all original columns and appends an \texttt{archived\_at} timestamp and an \texttt{archive\_reason} code. Preserving the full schema in the archive means that historical analysis and audit queries can run against archived data using the same column names and query patterns as the live table. Archival is triggered periodically to control the row count in the hot table and keep backend query latency within acceptable bounds during active fleet operations.

\section{Machine Learning Methodology}
\label{sec:ml}

The anomaly-detection pipeline uses Isolation Forest Model C as the primary production detector and Matrix Profile as a supporting subsequence detector. Model C receives context-aware features derived from fuel level, speed, distance, engine operation, driver behavior, route context, and fuel-economy baseline deviation. Matrix Profile is retained to support detection of gradual subsequence-level fuel-loss patterns. Model performance is evaluated using leave-one-trip-out validation so each held-out trip is scored by models trained and calibrated on separate trips.

\subsection{Anomaly Model and Environmental Constraints}
\label{sec:anomaly_model_constraints}

A fuel anomaly is formally defined as a telemetry row whose observed fuel-level change deviates from the expected fuel-level change under the current operating context. For row $t$, the normalized anomaly magnitude is represented as

\begin{equation}
A(t) =
\frac{\left|\Delta f_{\mathrm{obs}}(t) - \mathbb{E}[\Delta f \mid c(t)]\right|}
{\sigma(c(t))}
\label{eq:fuel_anomaly_model}
\end{equation}

\noindent
where $\Delta f_{\mathrm{obs}}(t)$ is the observed filtered fuel-level change, $\mathbb{E}[\Delta f \mid c(t)]$ is the expected fuel-level change under context $c(t)$, and $\sigma(c(t))$ is the expected spread of fuel change under that context. The context vector $c(t)$ includes speed, RPM, engine load, throttle, route, driver, and motion state. The system is therefore not a simple ``fuel dropped suddenly'' detector. It is a context-conditioned fuel-deviation detector that asks whether the observed fuel movement is plausible for the vehicle state in which it occurred.

The model is constrained first by the physical behaviour of the vehicle and tank. C1, conservation, states that fuel decreases through engine consumption or unauthorized removal, and increases only during refueling. C2, bounded consumption, states that a fuel drop cannot exceed what is physically plausible from engine work and load. C3, sender non-linearity, accounts for the float-arm sender being more unstable in certain tank-level bands. C4, slosh, recognizes that short fluctuations during motion can be caused by liquid movement rather than true consumption.

The model is also constrained by operational and environmental conditions in Philippine fleet use. O1 accounts for stop-and-go traffic and idling, where low distance travelled does not necessarily imply abnormal fuel loss. O2 accounts for route restrictions and truck-ban-related route differences that change expected consumption. O3 accounts for driver variability and per-driver fuel-economy differences. O4 accounts for the fact that refueling should occur during stationary or quiescent windows rather than during ordinary motion.

\begin{table}[!htbp]
\centering
\caption{Constraint-to-Mechanism Mapping for Fuel-Anomaly Detection}
\label{tab:constraint_mechanism_mapping}
\scriptsize
\setlength{\tabcolsep}{3pt}
\renewcommand{\arraystretch}{1.15}
\begin{tabularx}{\textwidth}{|p{0.22\textwidth}|X|X|}
\hline
\textbf{Constraint} & \textbf{System Mechanism} & \textbf{Purpose} \\
\hline
C1 Conservation & Direction-aware step rule and refuel quiescence window & Separate true fuel decreases and refueling events from impossible direction changes. \\
\hline
C2 Bounded consumption & \texttt{fuel\_drop\_per\_rpm}, \texttt{fuel\_drop\_per\_load}, and \texttt{expected\_consumption\_residual} features & Compare observed fuel loss against engine-work plausibility. \\
\hline
C3 Sender non-linearity & Median and EMA filter cascade tuned by tank behaviour & Reduce sender-band instability before model inference. \\
\hline
C4 Slosh & Median window and dead-band publisher & Suppress short fuel fluctuations caused by tank movement. \\
\hline
O1 Traffic and idling & \texttt{idle\_duration\_score} and \texttt{motion\_state\_code} & Distinguish stop-and-go operation from abnormal stationary loss. \\
\hline
O2 Route restrictions & \texttt{route\_type\_code} and route-segment performance features & Condition expected consumption on route type and segment behaviour. \\
\hline
O3 Driver variability & \texttt{driver\_kmpl\_deviation\_pct} and \texttt{driver\_behavior\_score} & Detect deviations from the driver's normal fuel-economy band. \\
\hline
O4 Refuel pattern & \texttt{time\_since\_refuel\_norm}, \texttt{post\_refuel\_flag}, and contextual refuel gate & Exempt plausible refueling and post-refuel stabilization windows. \\
\hline
\end{tabularx}
\normalsize
\end{table}

Table~\ref{tab:constraint_mechanism_mapping} links each physical or operational constraint to an implemented mechanism in the feature pipeline, filter cascade, or contextual gate. This makes the anomaly model auditable: every anomaly can be traced back to physical or operational constraints, not only to a black-box anomaly score.

\subsection{Machine Learning Architecture}

\begin{figure}[!htbp]
    \centering
    \includegraphics[width=0.95\textwidth]{ML_Anomaly_Detection_Architecture.png}
    \caption[Machine Learning Anomaly Detection Architecture]{Machine-learning anomaly-detection architecture organised into four planes. The \emph{data plane} captures OBD-II Mode~22 fuel, GPS, and HMI-state signals, applies the edge filter cascade (median + EMA, 5\,s HMAC-signed payload), ingests the row through \texttt{POST~/telemetry}, and derives the 33-dimensional Model~C feature vector plus a short fuel-window buffer. The \emph{primary inference plane} scores every row with Isolation Forest Model~C under the \texttt{if\_with\_margin} operating mode ($\text{score} \geq \tau + \delta$, with $\delta = 0.005$) and hands admitted candidates to the contextual gate. The \emph{audit-and-override inference plane} runs Matrix Profile with an $m = 4$ subsequence and reports a discord $z$-score for the same row; the score is persisted on the row as an audit signal, and additionally triggers an escalation arrow into the contextual gate when the discord exceeds the extreme-tail thresholds (\texttt{mp\_score}~$\geq 4.0$ with $\Delta\text{fuel} < 0$, or \texttt{mp\_score}~$\geq 2.0$ with $\Delta\text{fuel} \leq -2.0$\,L). The contextual gate is a five-stage cascade of one positive plausibility filter (\texttt{suspiciousContext}) and four negative suppressors (\texttt{parkedRefuel}, \texttt{postRefuel}, \texttt{loadSpikeSuppress}, \texttt{coldStartSuppress}); only rows that pass all five stages proceed to the persistence step, which sets \texttt{anomaly\_flag}, records the score and \texttt{model\_source}, tags the operational scenario from the 90\,s event window, and applies a 5-minute per-truck deduplication. The \emph{output plane} propagates the alert to the \texttt{alerts} table, drives the buzzer and banner on the driver HMI, and surfaces the event on the operator dashboard and \texttt{FleetMap} with the Matrix Profile audit column rendered as a second-opinion badge.}
    \label{fig:ml_anomaly_detection_architecture}
\end{figure}

Figure~\ref{fig:ml_anomaly_detection_architecture} illustrates the end-to-end anomaly-detection pipeline as a four-plane diagram: a data plane that converts raw sensor traffic into a persisted feature row, two inference planes that scored the row with complementary detectors, and an output plane that translates a passed candidate into an operator-facing alert. This section explains how a single 5-second telemetry row moves through each plane.

The data plane begins with the on-vehicle sensors --- the OBD-II Mode~22 fuel channel for the Toyota Hilux partner-fleet unit, the GPS NMEA stream, and the HMI trip-state signal --- and applies the fuel filter cascade documented in Section~\ref{sec:fuel_preprocessing}. The filtered payload is signed with a device-side HMAC, transmitted every 5\,s through the intelligent traffic-steering chain (LoRa~$\rightarrow$~LTE~$\rightarrow$~GSM~2G~$\rightarrow$~WiFi~$\rightarrow$~SPIFFS buffer), and ingested at the backend through the \texttt{POST~/telemetry} endpoint. The ingest step validates the row, resolves the trip, computes distance deltas through a Postgres RPC call, and derives the 33-feature Model~C vector that Section~\ref{sec:ml} enumerates. This is the point at which a row transitions from a raw sensor reading into a feature-carrying record and enters the machine-learning subsystem.

The primary inference plane operates row-by-row. The ingest layer invokes the ML service, which loads the pre-trained Isolation Forest Model~C pickle and scores the incoming feature vector. Under the \texttt{if\_with\_margin} operating mode the row is admitted as an anomaly candidate only when its anomaly score exceeds the per-context calibrated threshold $\tau$ by an empirical margin $\delta = 0.005$ carried over from the June~23 threshold-calibration pass. This margin is deliberately conservative because Isolation Forest scores are strongly bimodal on this corpus, so the added margin sacrifices a small fraction of borderline recall in exchange for a large reduction in false positives in production. Candidates that clear this test enter the contextual gate; candidates that do not are recorded in \texttt{telemetry\_logs} with an \texttt{anomaly\_score} but no \texttt{anomaly\_flag}, matching the ``IF quiet, no alert'' exit shown at the top-right of the primary plane in Figure~\ref{fig:ml_anomaly_detection_architecture}.

The audit-and-override inference plane operates on the same row in parallel, but at a subsequence granularity instead of a per-row granularity. The Matrix Profile subsystem maintains a rolling fuel-level buffer of the current trip's filtered fuel series and computes a $z$-score of the discord distance for the most recent length-$m = 4$ window against the calibration distribution derived from the clean reference trips in Section~\ref{sec:mp_math}. The output is a single scalar (\texttt{mp\_score}) and a companion boolean (\texttt{mp\_flag}) that is set whenever the score exceeds the calibrated tail threshold. Under the default operating mode this signal is treated as audit-only: it is persisted on the same \texttt{telemetry\_logs} row, and the dashboard uses it to display a second-opinion badge on every alert card. It does not, by itself, cause an alert to be raised. The exception is the \texttt{mpScreamOverride} escalation rule, which lets Matrix Profile trigger a primary alert when its discord score is unambiguously in the tail of the clean-trip distribution: either \texttt{mp\_score}~$\geq 4.0$ on a row with any negative $\Delta\text{fuel}$, or \texttt{mp\_score}~$\geq 2.0$ on a row whose per-row fuel loss is at least $2.0$\,L. This rule exists to catch slow gradual leaks whose per-row fuel movement is below every threshold in the primary plane but whose subsequence shape is a genuine discord. In Figure~\ref{fig:ml_anomaly_detection_architecture} the override is drawn as the arrow ascending from Matrix Profile into the contextual gate, and the rule remains gated by the same five-stage cascade that the primary path is subject to.

The contextual gate is the false-positive filter that both admission paths share. It is composed of one positive plausibility gate and four negative suppressors, and a row must satisfy the positive gate and clear all four suppressors before an alert is persisted. The positive gate, \texttt{suspiciousContext}, admits the row when at least one of four physical loss signatures is present: a per-row fuel delta strictly below $-2.0$\%, a fuel-per-kilometre rate above $0.6$\% of tank per km, a stationary idle-theft signature at speed below 5\,km/h with the engine on and fuel delta below $-1.0$\%, or a physically impossible driving-fuel increase of $+2$\% or more at speed above 5\,km/h. The four suppressors are then applied in sequence. \texttt{parkedRefuel} rejects the canonical fill-up signature (fuel delta above $+10$\% while parked). \texttt{postRefuel} rejects rows within five minutes of any $+5$\% fuel jump on the same trip and communication channel, because the capacitive fuel sender takes several minutes to settle after a fill. \texttt{loadSpikeSuppress} rejects the acceleration-slosh signature in which fuel briefly reads lower during a hard load transient (engine load $\geq 55$\% with a rise $\geq 25$\%, and a small negative fuel delta). \texttt{coldStartSuppress} rejects the sender-priming interval that follows an engine ignition transition. Rows blocked by any suppressor are logged and audited but do not generate operator-facing alerts, matching the ``suppressed, no alert'' exit shown at the top-right of the gate in Figure~\ref{fig:ml_anomaly_detection_architecture}. This separation of concerns --- an unsupervised statistical detector that produces candidates, followed by a domain-aware precision filter that suppresses the four families of sender artefacts that dominate the false-positive budget --- is what allows the deployed configuration to reach the 92.9\% precision reported in Chapter~\ref{ch:result_discuss}.

A row that passes the contextual gate reaches the persistence step. Persistence performs three writes. It updates the \texttt{telemetry\_logs} row with \texttt{anomaly\_flag}, \texttt{anomaly\_score}, and \texttt{model\_source} (\texttt{isolation\_forest} for a primary-path alert, \texttt{matrix\_profile} for an override-path alert). It runs a 90-second event-window scenario classifier that inspects the surrounding fuel drop shape and back-propagates the anomaly flag across the contiguous event so that the dashboard renders one alert per event instead of one alert per peak row. It then applies a 5-minute per-truck dedup window on the same \texttt{fuel\_anomaly} alert type before inserting a new row into the \texttt{alerts} table. A row that would pass the gate but whose combined score falls below the persistence confidence threshold ($\text{combined\_score} \leq 0.65$) or whose dedup check fails is recorded with its audit fields but is not inserted into the \texttt{alerts} table, matching the ``score $\leq 0.65$, dropped'' exit shown at the top-right of the persistence stage.

The output plane surfaces admitted alerts through three synchronized channels. The \texttt{alerts} row is the canonical record and carries the \texttt{fuel\_anomaly} type, the severity, and the human-readable scenario message. The driver HMI subscribes to the alert stream over the intelligent traffic-steering chain and drives a local buzzer plus banner with a \emph{Take Rest} or \emph{Snooze} affordance. The operator dashboard subscribes to the same stream and updates the alerts page and \texttt{FleetMap} panels in real time, additionally rendering the Matrix Profile audit column on each card so the operator can see at a glance whether the subsequence-shape detector agreed with the row-level detector. This is the diagram's bridge between methodology and evaluation: it explains why Chapter~\ref{ch:result_discuss} reports precision and recall from the model layer, false-positive rates from the gate layer, and closed-loop resolution statistics from the alert layer, and why all three are needed to justify a claim of production readiness.

\begin{figure}[!htbp]
\centering
\includegraphics[width=1\textwidth]{Fig_5_11_Anomaly_Pipeline.png}
\caption{Anomaly Modeling and Data-Mining Pipeline}
\label{fig:anomaly_modeling_data_mining_pipeline}
\end{figure}

Figure~\ref{fig:anomaly_modeling_data_mining_pipeline} summarizes the methodological path used in the implemented detector. Raw telemetry is first filtered, then converted into context features, screened by environmental constraints, scored by Isolation Forest Model~C, and finally translated into a constraint-violation signature that can be attached to an operator-facing alert.

\subsection{Dataset Composition}
The evaluation corpus merges 20{,}412 telemetry rows across 79 trips from two sources. The project-native partition contains the 634 rows captured by the prototype, on which the injected anomalies (six anomalous trips, 85 anomalous samples) were applied. The public partition is the 19{,}778-row publicly-derived corpus, retained without anomaly labels as the out-of-distribution clean-route domain.

% Keep the Chapter 5 table sequence continuous after the preceding captioned TikZ figure.
\setcounter{table}{7}
\begin{table}[!htbp]
\centering
\caption{Dataset Composition for Evaluation Corpus}
\label{tab:dataset_composition}
\scriptsize
\begin{tabular}{|p{0.55\textwidth}|p{0.22\textwidth}|}
\hline
\textbf{Dimension} & \textbf{Value} \\
\hline
Total telemetry rows & 20{,}412 \\
\hline
Total trips & 79 \\
\hline
Anomalous trips & 6 \\
\hline
Normal samples (clean) & 20{,}327 \\
\hline
Anomalous samples (labelled) & 85 \\
\hline
Project-native partition (prototype-captured) & 634 rows \\
\hline
Public partition (publicly-derived clean routes) & 19{,}778 rows \\
\hline
Injected labelled rows & 169 rows \\
\hline
Headline evaluation subset & 169 rows \\
\hline
\end{tabular}
\normalsize
\end{table}

The headline performance is computed against the project-native partition, which isolates the injected anomalies on telemetry produced by the prototype sensor stack. The public partition is reported separately as a stress test, since pooling it would dilute the precision--recall denominators with negatives from a different driving domain.

Table~\ref{tab:dataset_composition} therefore defines the denominator for every SO4 claim. The project-native partition answers whether the prototype's own telemetry and injected labels can be detected, while the public partition answers whether the detector remains quiet on a larger clean-route domain. Keeping these roles separate avoids overstating performance through class imbalance.

\subsection{Anomaly Injection Methodology}
\label{sec:anomaly_injection}

\subsubsection{Why Injection Is Necessary}
Real fuel theft is unsafe to stage under partner-fleet conditions (siphoning gasoline is a fire hazard and the partner operator did not authorise the risk), and confirmed theft events at SGSA reach the fleet operator only days after the fact through odometer-versus-invoice reconciliation, so the raw telemetry rows corresponding to those events are unlabelled at the row level. Controlled synthetic injection is therefore both required and standard: Turan et al.\ (2024), Aquize et al.\ (2017), and Barbado and Corcho (2022) all evaluate against injected anomalies for the same reason. This section defines the injection model formally and then documents, for each scenario, the physical story, the falsified fuel curve with all symbols named, a worked numeric example, and a combustion-based physics check demonstrating that the injected rate is thermodynamically distinguishable from real burn.

\subsubsection{Formal Injection Tuple}

For a real telemetry row $r_t = \bigl(\mathrm{truck\_id},\, \mathrm{trip\_id},\, \tau_t,\, \mathrm{gps}_t,\, F_t,\, s_t,\, \omega_t,\, L_t,\, \theta_t\bigr)$, the injected row is
\begin{equation}
r'_t = r_t \;\oplus\; \bigl(\delta F_t,\, \delta s_t,\, \delta \omega_t,\, \delta L_t,\, \delta \theta_t\bigr) \cdot \mathbb{1}\bigl[\,t \in W_{\mathrm{scenario}}\,\bigr]
\label{eq:injection_model}
\end{equation}
where $F_t$ is the fuel percentage, $s_t$ is the vehicle speed in km/h, $\omega_t$ is the engine RPM, $L_t$ is the engine load in \%, and $\theta_t$ is the throttle position in \%. The tuple $(\delta F_t,\, \delta s_t,\, \delta \omega_t,\, \delta L_t,\, \delta \theta_t)$ contains the scenario-defined override values. The window $W_{\mathrm{scenario}} \subset \mathbb{Z}$ is a contiguous set of sample indices during which the injection is active; the indicator $\mathbb{1}[\cdot]$ turns injection on inside the window and off outside. The operator $\oplus$ denotes field-wise \emph{overwrite}, not addition~---~inside the window the tuple's five fields replace their counterparts in $r_t$, and outside the window $r'_t = r_t$. The ground-truth label $y_t \in \{0, 1\}$ is written alongside $r'_t$ as $y_t = 1 \Leftrightarrow t \in W_{\mathrm{scenario}}$; this is the reference against which precision and recall are computed. Identifiers, timestamps, GPS, and odometer are preserved so that the injection is grounded in real trip context; the machine-learning pipeline receives $r'_t$ with no knowledge of the injection window, the scenario, or the label.

\subsubsection{Universal Symbol Dictionary for Per-Scenario Curves}

The three scenarios below share a common symbol set. Let $t$ denote the sample index with fixed cadence $\Delta_s = 5$~s (telemetry publish cadence, Section~\ref{sec:fuel_notation}), $t_0$ the sample at which the injection window \emph{begins}, and $\Delta t_{\mathrm{win}}$ the window \emph{duration in seconds}. The set $W = \{t : t_0 \leq t \leq t_0 + \Delta t_{\mathrm{win}}\}$ is the sample-index window. Let $F_0$ (in \%) be the fuel percentage at $t_0$, $F_{\mathrm{end}} = F_0 - \beta_k \cdot \Delta t_{\mathrm{win}}$ the fuel percentage at $t_0 + \Delta t_{\mathrm{win}}$, and $\beta_k$ (in \%/s) the drain-rate slope of scenario $k$. The total injected magnitude is $\Delta F_{\mathrm{total}} = F_0 - F_{\mathrm{end}}$. The companion overrides $\delta s_t$, $\delta \omega_t$, $\delta L_t$, $\delta \theta_t$ are the scenario-consistent speed, RPM, load, and throttle values held constant across $W$ so that the injection tuple is physically self-consistent.

\subsubsection{Scenario~1: Stationary Siphon (Overnight-Theft Signature)}
\label{sec:inject_scenario1}

\textbf{Physical story.} A truck is parked at a warehouse gate overnight with the ignition off. An operator opens the fuel cap and runs a hand-pump siphon into a jerry can. No headlights come on, no engine noise, no motion~---~only the tank level falls. This signature corresponds to the SGSA NAIA-staging-area incident report of 2025-11-14.

\textbf{Falsified fuel curve.}
\begin{equation}
F_t = F_0 - \beta_1 \cdot (t - t_0) \cdot \Delta_s \quad \text{for all } t \in W, \qquad \beta_1 = 0.56\,\%/\text{s}, \quad \Delta t_{\mathrm{win}} = 25\text{ s}, \quad \Delta F_{\mathrm{total}} = 14\,\%
\label{eq:injection_s1}
\end{equation}
with companion overrides $\delta s_t = 0$~km/h, $\delta \omega_t = 0$~rev/min (engine off), $\delta L_t = 0\%$, $\delta \theta_t = 0\%$ for all $t \in W$.

\textbf{Worked example.} Take $F_0 = 78\%$ (matches the \texttt{simulate\_scenarios.py} \texttt{N5} scenario \texttt{start\_fuel}). Sample-by-sample the injected trace is $F(t_0) = 78.0\%$, $F(t_0 + 5) = 78.0 - 0.56 \cdot 5 = 75.2\%$, $F(t_0 + 10) = 78.0 - 0.56 \cdot 10 = 72.4\%$, and $F(t_0 + 25) = 78.0 - 0.56 \cdot 25 = 64.0\%$. After the window the trace is frozen at $F_{\mathrm{end}} = 64\%$ because fuel does not spontaneously reappear.

\textbf{Physics check.} The idle-combustion upper bound on the fuel-loss rate for a 2.4~L diesel is $\dot{V}_{\mathrm{idle}} \approx 0.5$~L/hr (SAE~J1349), which on the 80~L tank produces a percentage rate of $0.5/(80 \cdot 60) = 0.00010$~\%/s $= 0.00029$~\%/s counting engine-off electrical draw for auxiliaries. The injected rate is $\beta_1 / 0.00029 \approx 1{,}900\times$ the maximum plausible idle-combustion rate. Because the engine is off in the companion overrides, there is no combustion at all: the ML detector observes fuel falling under $L_t = 0$ and $s_t = 0$, a state that the trained context features assign a vanishingly small probability. The primary detection route in Model~C is the \texttt{load\_per\_speed\_norm} feature (row~23 in Table~\ref{tab:if_feature_set_variants}). Supporting evidence for this signature appears in Aquize et al.\ (2017) and Beteto et al.\ (2022).

\subsubsection{Scenario~2: Rapid Mid-Drive Drop (Line-Tap During Dispatch)}
\label{sec:inject_scenario2}

\textbf{Physical story.} A truck is moving at low speed inside a depot yard when a second vehicle rolls alongside; an external pump is briefly attached to the fuel-return line. The driver observes nothing on the dashboard gauge because factory sender damping hides fast changes, but the OBD-II PID stream captures the fall in real time. This is the shortest and highest-slope signature we inject and corresponds to the SGSA Cubao depot-yard report of 2026-02-08.

\textbf{Falsified fuel curve.}
\begin{equation}
F_t = F_0 - \beta_2 \cdot (t - t_0) \cdot \Delta_s \quad \text{for all } t \in W, \qquad \beta_2 = 0.62\,\%/\text{s}, \quad \Delta t_{\mathrm{win}} = 25\text{ s}, \quad \Delta F_{\mathrm{total}} = 15.5\,\%
\label{eq:injection_s2}
\end{equation}
with companion overrides $\delta s_t = 45$~km/h, $\delta \omega_t \approx 1{,}800$~rev/min, $\delta L_t = 55\%$, $\delta \theta_t \approx 40\%$.

\textbf{Worked example.} Take $F_0 = 68\%$ (matches the \texttt{simulate\_scenarios.py} \texttt{a1} scenario \texttt{start\_fuel}). The injected trace is $F(t_0) = 68.0\%$, $F(t_0 + 5) = 68.0 - 0.62 \cdot 5 = 64.9\%$, $F(t_0 + 15) = 58.7\%$, and $F(t_0 + 25) = 52.5\%$. Each 5-s step drops approximately 3.1\%, which is above the $-2.5\%$ direction-aware gate threshold from Equation~\ref{eq:direction_aware}, so each step is published by the filter cascade and reaches the ML layer intact.

\textbf{Physics check.} At $s_t = 45$~km/h under $L_t = 55\%$ engine load, aggressive burn is approximately 8~L per 100~km, or $3.6$~L/hr, which on the 80~L tank gives $0.075$~\%/min $= 0.00125$~\%/s. The injected rate is $\beta_2 / 0.00125 \approx 500\times$ the maximum plausible driving-burn rate. Even a stuck-open injector is bounded by the fuel-pump flow rate, which is well below $\beta_2$; the injected drop is thermodynamically impossible from combustion alone. The primary detection route in Model~C is a spike in the \texttt{kmpl\_deviation\_pct} feature (row~31 in Table~\ref{tab:if_feature_set_variants}) because fuel-per-kilometre explodes over the 25~s window. Supporting evidence for this signature appears in Turan et al.\ (2024) and Barbado and Corcho (2022).

\subsubsection{Scenario~3: Gradual Leak (Long-Haul Rest-Stop Drilling)}
\label{sec:inject_scenario3}

\textbf{Physical story.} A long-haul truck stops at a highway rest area. A 1--2~mm hole is drilled in the underside of the tank; fuel drains slowly over approximately 40~s. The driver may not notice for hours because the gauge integrates too slowly. This is the hardest case for a point-anomaly detector because each individual 5-s step lies inside the sensor-noise envelope of the filter cascade; only the multi-sample subsequence pattern reveals the drop, which is why the two-window Matrix Profile in Section~\ref{sec:mp_math} is retained as a companion detector.

\textbf{Falsified fuel curve.}
\begin{equation}
F_t = F_0 - \beta_3 \cdot (t - t_0) \cdot \Delta_s \quad \text{for all } t \in W, \qquad \beta_3 = 0.135\,\%/\text{s}, \quad \Delta t_{\mathrm{win}} = 40\text{ s}, \quad \Delta F_{\mathrm{total}} = 5.4\,\%
\label{eq:injection_s3}
\end{equation}
with companion overrides $\delta s_t = 35$~km/h, $\delta \omega_t \approx 1{,}500$~rev/min, $\delta L_t = 40\%$, $\delta \theta_t \approx 25\%$.

\textbf{Worked example.} Take $F_0 = 51\%$. The injected trace is $F(t_0) = 51.00\%$, $F(t_0 + 10) = 51.00 - 0.135 \cdot 10 = 49.65\%$, $F(t_0 + 20) = 48.30\%$, and $F(t_0 + 40) = 45.60\%$. Each 5-s step drops approximately $0.68\%$~---~below the $-2.5\%$ direction-aware gate threshold and below the moving dead-band of $0.5\%$, so the filter cascade admits each step as normal. Over the 40~s window, however, the aggregate 5.4\% drop is anomalous for the $m = 12$-sample Matrix Profile subsequence.

\textbf{Physics check.} At $s_t = 35$~km/h under $L_t = 40\%$ engine load, cruise burn is approximately 7.5~L per 100~km $= 2.6$~L/hr $= 0.055$~\%/min $= 0.00092$~\%/s. The injected rate is $\beta_3 / 0.00092 \approx 135\times$ the plausible cruise-burn maximum. The factor is deliberately smaller than for Scenarios~1--2 because the entire point of Scenario~3 is that a point detector cannot see it: the necessity of the Matrix Profile layer is defended by this case. Supporting evidence for this signature appears in Fortela et al.\ (2023) and Garc\'{\i}a Vega (2022).

\subsubsection{At-a-Glance Layer Coverage}

Table~\ref{tab:injection_layer_coverage} maps each scenario to the pipeline layer whose behaviour it specifically stresses; dropping any layer causes at least one scenario to slip.

\begin{table}[!htbp]
\centering
\caption{Injected Scenarios and the Pipeline Layers They Stress}
\label{tab:injection_layer_coverage}
\scriptsize
\begin{tabular}{|p{0.03\textwidth}|p{0.08\textwidth}|p{0.09\textwidth}|p{0.10\textwidth}|p{0.10\textwidth}|p{0.38\textwidth}|}
\hline
\textbf{\#} & $\beta$ (\%/s) & $\Delta t_{\mathrm{win}}$ (s) & $\Delta F_{\mathrm{total}}$ (\%) & \textbf{$\times$ burn max} & \textbf{Detector that must catch it} \\
\hline
1 & 0.56 & 25 & 14.0 & $\approx 1900\times$ & Isolation Forest via \texttt{load\_per\_speed\_norm} (engine off + fuel falling) \\
\hline
2 & 0.62 & 25 & 15.5 & $\approx 500\times$ & Isolation Forest via \texttt{kmpl\_deviation\_pct} (fuel/km explodes) \\
\hline
3 & 0.135 & 40 & 5.4 & $\approx 135\times$ & Matrix Profile $m = 12$ discord (IF is architecturally blind) \\
\hline
\end{tabular}
\normalsize
\end{table}

\subsubsection{Scenario Coverage in Chapter~6 Evaluation}
Table~\ref{tab:anomaly_scenario_coverage} maps the three formal scenarios above onto the broader operational classes used in the Chapter~\ref{ch:result_discuss} results, adding the sensor-spike and heavy-legitimate-consumption cases required to verify the filter cascade and the false-positive discipline of the contextual features.

\begin{table}[!htbp]
\centering
\caption{Anomaly Scenario Coverage}
\label{tab:anomaly_scenario_coverage}
\scriptsize
\begin{tabular}{|p{0.20\textwidth}|p{0.32\textwidth}|p{0.22\textwidth}|p{0.12\textwidth}|}
\hline
\textbf{Scenario} & \textbf{Description} & \textbf{Primary detector} & \textbf{Coverage} \\
\hline
Sudden fuel theft & Single-step drop of at least 10\% & Isolation Forest & Yes \\
\hline
Gradual siphoning & Slow steady decrease & Matrix Profile & Yes \\
\hline
Continuous leak & Linear decline at speed = 0 & Matrix Profile & Yes \\
\hline
Sensor spike & Brief out-of-range raw value & Median filter (pre-ML) & Yes \\
\hline
Heavy legitimate consumption & Loaded uphill drive & Both models correctly classify as normal & Yes \\
\hline
\end{tabular}
\normalsize
\end{table}

\noindent
Table~\ref{tab:anomaly_scenario_coverage} documents why the evaluation is not limited to one obvious theft pattern. Sudden fuel theft, gradual siphoning, continuous leakage, sensor spikes, and heavy legitimate consumption stress different parts of the system: the filter should reject sensor spikes, the ML model should detect abnormal fuel loss, and the contextual features should avoid flagging legitimate heavy consumption as theft. The scenario coverage becomes the basis for the per-class interpretation in Chapter~6.

\subsection{Feature Engineering}

The machine-learning methodology used three Isolation Forest variants so the effect of additional context could be tested step by step.\glsadd{not:fi}\glsadd{not:deltafi}\glsadd{not:di}\glsadd{not:vi}\glsadd{not:kmpli} Model A is the baseline context-aware model and uses fuel behavior, motion state, engine operation, and refuel-window indicators. Model B extends Model A by adding speed and RPM band-share features that summarize how intensely the vehicle was being driven. Model C extends Model B by adding fuel-economy deviation features for the fleet vehicle, driver, and route context; this is the final production model because it explains whether a fuel change is abnormal relative to the operating situation rather than only relative to the immediately preceding samples. Matrix Profile was evaluated as the time-series shape detector and as part of the hybrid comparison, but the Chapter~6 results show that the Model C Isolation Forest provided the best balance of precision, recall, and false-positive control.

The feature names in Table~\ref{tab:if_feature_set_variants} are kept in their backend form for reproducibility, while the ``What it captures'' column gives the plain-language role of each feature. This avoids treating source-code identifiers as the argument of the thesis while still allowing the results to be traced back to the implemented system.

\begingroup
\begin{center}
\tiny
\setlength{\tabcolsep}{1.5pt}
\renewcommand{\arraystretch}{1.12}
\begin{longtable}{|p{0.035\textwidth}|p{0.19\textwidth}|p{0.19\textwidth}|p{0.30\textwidth}|p{0.035\textwidth}|p{0.035\textwidth}|p{0.035\textwidth}|}
\caption{Feature Set Across Isolation Forest Variants (A / B / C)}
\label{tab:if_feature_set_variants} \\
\hline
\textbf{\#} & \textbf{Feature} & \textbf{Source signal} & \textbf{What it captures} & \textbf{A} & \textbf{B} & \textbf{C} \\
\hline
\endfirsthead
\hline
\textbf{\#} & \textbf{Feature} & \textbf{Source signal} & \textbf{What it captures} & \textbf{A} & \textbf{B} & \textbf{C} \\
\hline
\endhead
1 & \texttt{fuel\_level\_filtered} & OBD-II PID 21/29 & Baseline tank level after filter cascade & Y & Y & Y \\
\hline
2 & \texttt{speed\_filtered} & OBD-II PID 0x0D + GPS fallback & Smoothed vehicle speed & Y & Y & Y \\
\hline
3 & \texttt{fuel\_delta} & \texttt{fuel\_level\_filtered} & Per-row consumption increment & Y & Y & Y \\
\hline
4 & \texttt{odometer\_delta} & OBD-II odometer / integrated speed & Distance moved per sample & Y & Y & Y \\
\hline
5 & \texttt{fuel\_per\_km} & \texttt{fuel\_delta} / \texttt{odometer\_delta} & Efficiency rate and siphoning-under-load signal & Y & Y & Y \\
\hline
6 & \texttt{fuel\_rate\_per\_hour} & \texttt{fuel\_delta} / dt & Continuous-leak signature & Y & Y & Y \\
\hline
7 & \texttt{distance\_per\_fuel} & 1 / \texttt{fuel\_per\_km} & Light-load / coasting regime & Y & Y & Y \\
\hline
8 & \texttt{route\_segment\_performance} & Per-segment mean vs trip mean & Route-specific consumption shift & Y & Y & Y \\
\hline
9 & \texttt{driver\_behavior\_score} & Acceleration, variability, and idle aggregate & Long-tail driving-style signal & Y & Y & Y \\
\hline
10 & \texttt{engine\_rpm\_normalized} & OBD-II PID 0x0C & Engine work intensity & Y & Y & Y \\
\hline
11 & \texttt{engine\_load\_pct\_norm} & OBD-II PID 0x04 & ECU-reported load fraction & Y & Y & Y \\
\hline
12 & \texttt{throttle\_pct\_norm} & OBD-II PID 0x11 & Driver demand & Y & Y & Y \\
\hline
13 & \texttt{fuel\_drop\_while\_stationary} & \texttt{fuel\_delta} + speed & Siphoning signature while parked & Y & Y & Y \\
\hline
14 & \texttt{fuel\_drop\_with\_engine\_off} & \texttt{fuel\_delta} + engine status & Theft-at-rest signature & Y & Y & Y \\
\hline
15 & \texttt{idle\_duration\_score} & Rolling idle flag sum & Sustained-idle context & Y & Y & Y \\
\hline
16 & \texttt{fuel\_drop\_per\_rpm} & Fuel use / RPM & Disagreement between fuel use and engine work & Y & Y & Y \\
\hline
17 & \texttt{fuel\_drop\_per\_load} & Fuel use / engine load & Disagreement between fuel use and reported load & Y & Y & Y \\
\hline
18 & \texttt{motion\_state\_code} & Speed + engine status & Motion-regime category & Y & Y & Y \\
\hline
19 & \texttt{time\_since\_refuel\_norm} & Refuel-event detector & Refuel-window exemption signal & Y & Y & Y \\
\hline
20 & \texttt{post\_refuel\_flag} & Refuel-event detector & Post-refuel exemption flag & Y & Y & Y \\
\hline
21 & \texttt{delta\_engine\_load\_norm} & Engine-load delta & Rate of change of load & Y & Y & Y \\
\hline
22 & \texttt{delta\_throttle\_norm} & Throttle delta & Rate of change of driver input & Y & Y & Y \\
\hline
23 & \texttt{load\_per\_speed\_norm} & Load / speed & High load while parked or slow & Y & Y & Y \\
\hline
24 & \texttt{expected\_consumption\_residual} & Expected vs observed fuel use & Consumption residual against baseline & Y & Y & Y \\
\hline
25 & \texttt{rpm\_high\_share} & Engine RPM & Rolling fraction with RPM $>$ 2000 & - & Y & Y \\
\hline
26 & \texttt{rpm\_red\_share} & Engine RPM & Rolling fraction with RPM $>$ 3500 & - & Y & Y \\
\hline
27 & \texttt{rpm\_orange\_share} & Engine RPM & Rolling fraction with RPM 2500--3500 & - & Y & Y \\
\hline
28 & \texttt{rpm\_yellow\_share} & Engine RPM & Rolling fraction with RPM 1500--2500 & - & Y & Y \\
\hline
29 & \texttt{speed\_over\_90\_share} & Vehicle speed & Rolling fraction with speed $>$ 90 km/h & - & Y & Y \\
\hline
30 & \texttt{speed\_over\_120\_share} & Vehicle speed & Rolling fraction with speed $>$ 120 km/h & - & Y & Y \\
\hline
31 & \texttt{kmpl\_deviation\_pct} & Per-asset fuel-economy baseline & Vehicle-specific fuel-economy deviation & - & - & Y \\
\hline
32 & \texttt{driver\_kmpl\_deviation\_pct} & Per-driver fuel-economy baseline & Driver-specific fuel-economy deviation & - & - & Y \\
\hline
33 & \texttt{route\_type\_code} & Trip metadata + GPS context & Route-context disambiguator for KMPL\glsadd{KMPL} features & - & - & Y \\
\hline
\end{longtable}
\normalsize
\end{center}
\endgroup

The 33 features span three progressively extended groups. Model A contains 24 context-aware baseline features covering fuel consumption rates, OBD-II operational signals, motion state, and refuel-event flags. Model B extends Model A to 30 features by adding driving-zone band-share features derived from RPM and speed distributions, included for ablation comparison only. Model C extends Model B to 33 features by adding per-asset and per-driver fuel-economy baseline deviation features and a route-type code disambiguator, and is the production configuration used in Chapter~6.

\subsection{Data-Mining Layer}
\label{sec:data_mining_layer}

Isolation Forest Model~C is the row-level anomaly scorer used by the deployed ML pipeline. A separate data-mining layer was included to discover operating regimes and constraint-violation signatures from the telemetry data. This layer does not introduce a new production classifier; it explains how flagged rows relate to the physical and operational constraints defined in Section~\ref{sec:anomaly_model_constraints}.

\subsubsection{Constraint-Violation Profiling}

Every row flagged by the model can be annotated with constraint signatures such as C1, C2, C4, O1, O3, and O4. The annotation rules are deterministic and are applied after scoring so that the alert can be interpreted. C1 fires when a fuel drop is observed while parked. C2 fires when the per-step fuel drop exceeds plausible engine consumption. C4 fires when the raw-versus-filtered gap indicates slosh suppression. O1 fires when fuel drop occurs inside idle windows. O3 fires when the driver KMPL deviation exceeds the baseline band. O4 fires when the row is inside a post-refuel quiescence window.

\subsubsection{Archetype Discovery}

K-Means clustering with $K = 4$ was used only as an unsupervised data-mining layer to discover natural operating regimes inside the controlled-injection window. The clustering used seven $z$-standardized features: \texttt{fuel\_delta\_\allowbreak filtered}, \texttt{speed\_\allowbreak filtered}, \texttt{engine\_load\_pct\_\allowbreak norm}, \texttt{idle\_duration\_\allowbreak score}, \texttt{motion\_state\_\allowbreak code}, \texttt{throttle\_pct\_\allowbreak norm}, and \texttt{engine\_rpm\_\allowbreak normalized}. Standardization allows speed, fuel delta, load, and state codes to contribute on comparable scales.

The value $K = 4$ was selected because it separated four interpretable operating regimes in the controlled-injection window: high-speed cruise, mid-speed cruise, gradual-decline behaviour, and raw-slosh artifact behaviour. K-Means is not the deployed anomaly detector and does not replace Model~C. Its role is to support interpretation by showing which operating regime a flagged row belongs to and whether that regime is consistent with the constraint signature attached to the alert.

\subsection{Isolation Forest Training Methodology}
\label{sec:if_math}

The Isolation Forest is trained as an unsupervised outlier detector because the normal operating data are more abundant than confirmed abnormal fuel-loss events. Rather than modelling the density of the normal class, the forest exploits the observation that anomalies are easier to \emph{isolate} by random axis-parallel splits: they live in sparse regions of feature space and are separated from the rest of the sample in only a few cuts, whereas normal rows are surrounded by neighbours and require many cuts to isolate (Liu, Ting, and Zhou, 2008).

\subsubsection{Score Equation and Full Symbol Definition}

The Isolation Forest anomaly score is
\begin{equation}
S_{\mathrm{IF}}(x) = 2^{\,-\bar{h}(x) / c(n)}
\label{eq:if_score}
\end{equation}
where $x \in \mathbb{R}^{d}$ with $d = 33$ is one telemetry row after standardisation, and the mean isolation path length over the forest is
\begin{equation}
\bar{h}(x) = \frac{1}{T} \sum_{i=1}^{T} h_i(x)
\label{eq:if_mean_path}
\end{equation}
with $T = 300$ the number of trees in the forest (the \texttt{n\_estimators} hyperparameter) and $h_i(x) \in \mathbb{N}$ the number of edges traversed in tree $i$ from the root to the leaf that isolates $x$. The normaliser $c(n)$ is the expected path length of an unsuccessful search in a random binary search tree of size $n$ (Preiss, 1999):
\begin{equation}
c(n) = 2\cdot H(n-1) - \frac{2\,(n-1)}{n}, \qquad H(k) = \sum_{j=1}^{k} \frac{1}{j} \approx \ln(k) + \gamma
\label{eq:c_of_n}
\end{equation}
where $n = 256$ is the subsampling size per tree (the \texttt{max\_samples} hyperparameter, denoted $\psi$ by Liu et al.), $H(k)$ is the $k$-th harmonic number, and $\gamma \approx 0.5772$ is the Euler-Mascheroni constant. The score is bounded $S_{\mathrm{IF}}(x) \in (0, 1)$ with the following behaviour: $\bar{h}(x) \to 0 \Rightarrow S_{\mathrm{IF}} \to 1$ (extremely easy to isolate; strongly anomalous), $\bar{h}(x) \to c(n) \Rightarrow S_{\mathrm{IF}} \to 0.5$ (typical difficulty), and $\bar{h}(x) \to \infty \Rightarrow S_{\mathrm{IF}} \to 0$ (deep in a dense cluster; strongly normal). The bounded range and fixed reference at $0.5$ are what allow the threshold in Section~\ref{sec:contextual_gate_math} to be expressed as a percentile of clean-training scores rather than as an absolute cutoff.

\subsubsection{Hyperparameter Justification}
Table~\ref{tab:if_hyperparameter_justification} lists every hyperparameter with its statistical or literature-based justification. Each value was verified by sweep against the leave-one-trip-out (LOTO) folds rather than adopted from defaults without validation.

\begin{table}[!htbp]
\centering
\caption{Isolation Forest Hyperparameters and Justification}
\label{tab:if_hyperparameter_justification}
\scriptsize
\setlength{\tabcolsep}{3pt}
\renewcommand{\arraystretch}{1.20}
\begin{tabularx}{\textwidth}{|p{0.18\textwidth}|p{0.14\textwidth}|X|}
\hline
\textbf{Symbol / setting} & \textbf{Value} & \textbf{Statistical or literature basis} \\
\hline
$T$ (\texttt{n\_estimators}) & 300 & Liu et al.\ (2008) show path-length variance converges as $\mathcal{O}(1/T)$. A LOTO sweep over $T \in \{50, 100, 200, 300, 500, 1000\}$ produced a variance elbow near $T \approx 250$ (change less than 0.5~pp between $T = 300$ and $T = 1000$). $T = 300$ is the smallest tree count past the elbow. \\
\hline
$\psi$ (\texttt{max\_samples}) & 256 & Liu et al.\ (2008, \S4.2) prove that $\psi = 256$ saturates $c(\psi)$ for skewed data. Re-verified on the 634-row native partition: precision peaks at $\psi = 256$, drops 3~pp at $\psi = 128$ (undersampled trees miss context modes), and drops 1~pp at $\psi = 512$ (overfitting normal cluster). \\
\hline
$\nu$ (\texttt{contamination}) & 0.015 (1.5\%) & Three independent priors agree: the SGSA operational report of 3 confirmed theft events across 4 assets in 12 months gives approximately 0.75\%/asset/yr; Barbado and Corcho (2022) assume 1--3\%; Turan et al.\ (2024) use 2\%. A LOTO sweep over $\nu \in \{0.005, 0.01, 0.015, 0.02, 0.03\}$ maximised precision at $\nu = 0.015$. \\
\hline
Splitting rule & Random axis-aligned & Liu et al.\ (2008) canonical. Extended Isolation Forest (Hariri et al., 2019) with random-direction splits was tested and gained 0.4~pp precision but at $4\times$ training time and with loss of per-feature attribution; the axis-aligned rule was retained. \\
\hline
Threshold quantile & $q_{0.995}$ per context & Extreme-value literature (Coles, 2001) supports the 99.5-percentile as a standard tail threshold for rare events. Per-context computation is required because the clean-training $S_{\mathrm{IF}}$ distribution is bimodal across $s_t$ bins (idle vs.\ highway); a single global cutoff over-fires on idle and under-fires on highway. Pooled cutoff yielded FPR $= 4.1\%$ versus per-context $1.09\%$ on the 6 anomalous trips. \\
\hline
Operational margin $\delta$ & 0.003 & From the 2026-06-23 precision-margin calibration study on the project-native partition. Sweeping $\delta \in \{0, 0.001, 0.003, 0.005, 0.01\}$ produced a precision plateau at $\delta = 0.003$; $\delta > 0.005$ collapsed recall by more than 8~pp, and $\delta = 0$ admitted borderline scores that were 30\% false-positive. \\
\hline
\end{tabularx}
\normalsize
\end{table}

\subsection{Matrix Profile Methodology}
\label{sec:mp_math}

The Matrix Profile is retained as a subsequence-shape detector because the point-based Isolation Forest is architecturally blind to gradual leaks whose individual per-row deltas lie inside the filter-cascade noise envelope (Section~\ref{sec:inject_scenario3}). The two detectors are combined as an audit pair: agreement raises confidence, disagreement is the operational alert signal.

\subsubsection{Matrix Profile Equation}
For a $z$-normalised univariate fuel time series $\hat{T} \in \mathbb{R}^{N}$ of length $N$ derived from \texttt{fuel\_level\_filtered}, the Matrix Profile value at index $i$ is
\begin{equation}
P_i = \min_{\substack{j \neq i \\ |i-j| \geq m}} \bigl\| \hat{T}_{i:i+m} - \hat{T}_{j:j+m} \bigr\|_{2}
\label{eq:matrix_profile}
\end{equation}
where $m \in \mathbb{N}$ is the subsequence length, $\hat{T}_{k:k+m} \in \mathbb{R}^{m}$ is the length-$m$ subsequence starting at index $k$, and $\|\cdot\|_{2}$ is the Euclidean distance. The exclusion constraint $|i - j| \geq m$ prevents trivial self-matches from overlapping subsequences. The value $P_i$ is the distance from the subsequence at $i$ to its nearest non-trivial neighbour: small $P_i$ means the subsequence repeats often, large $P_i$ means the subsequence is a \emph{discord}~---~a candidate anomaly (Yeh et al., 2016).

The pipeline runs two windows in parallel. The primary window $m_1 = 5$ samples corresponds to a $5 \cdot \Delta_s = 25$~s subsequence, matched to the Scenario~2 rapid-drop signature. The long window $m_2 = 12$ samples corresponds to $12 \cdot \Delta_s = 60$~s, matched to the Scenario~3 gradual-leak signature. Sweeping $m \in \{3, 5, 8, 12, 20, 30\}$ demonstrated that Scenarios~2 and~3 have disjoint optimal window lengths, justifying the two-window design over any single-$m$ compromise.

\subsubsection{Robust Discord Threshold}
The discord threshold is applied on a Median Absolute Deviation (MAD) $z$-score rather than the standard-deviation $z$-score, because the very outliers we wish to detect would inflate the standard deviation and mask themselves:
\begin{equation}
z_i = \frac{P_i - \operatorname{median}(P)}{\operatorname{MAD}(P)}, \qquad \operatorname{MAD}(P) = \operatorname{median}\bigl(|P_i - \operatorname{median}(P)|\bigr)
\label{eq:mp_zscore}
\end{equation}
where $\operatorname{median}(P)$ is the sample median over all length-$m$ subsequences in the trip and $\operatorname{MAD}(P)$ is the median absolute deviation, a robust scale estimator that replaces the standard deviation (Rousseeuw and Croux, 1993). A row is flagged if $z_i > 2.5$, which under a Gaussian noise assumption corresponds to approximately $3.7\sigma$ and to an expected false-positive rate below one in four thousand. On the 27-trip corpus with a median trip length of 60~min = 720 subsequences at $m = 5$, this gives an expected fewer than 0.2 false positives per trip, consistent with the reported 1.09\% row-level FPR in Chapter~\ref{ch:result_discuss}. Alternative thresholds $z > 2.0$ (3.2 false positives per trip) and $z > 3.0$ (missed 2 of 6 anomalous trips) were rejected as Pareto-dominated.

In addition to the discord threshold, the Matrix Profile detector applies three deterministic post-discord gates: a magnitude gate ($|\Delta F| \geq 1\%$ absolute fuel change), a direction gate (downward-only, because leaks and thefts cannot cause fuel to rise), and a refuel-exemption gate (suppressing flags inside the post-refuel quiescence window from Stage~5 of Section~\ref{sec:fuel_preprocessing}). Each gate is calibrated against the held-out clean project-native trips at each LOTO fold.

\subsection{Contextual Gate and Alert Decision Rule}
\label{sec:contextual_gate_math}

The final alert decision combines the Isolation Forest score, the per-context threshold, and a conjunction of physical-plausibility gates:
\begin{equation}
\mathrm{alert}_t = \mathbb{1}\bigl[\, S_{\mathrm{IF}}(x_t) > \tau_{\mathrm{ctx}} + \delta \,\bigr] \;\wedge\; G_{\mathrm{context}}(x_t)
\label{eq:alert_rule}
\end{equation}
where $\tau_{\mathrm{ctx}} = q_{0.995}\bigl(S_{\mathrm{IF}} \mid \mathrm{context}(x_t)\bigr)$ is the 99.5-percentile of clean-training scores within the same context bin (Table~\ref{tab:if_hyperparameter_justification}), $\delta = 0.003$ is the operational margin, and $\mathrm{context}(x_t) \in \{\mathrm{idle}, \mathrm{low\_speed}, \mathrm{mid\_speed}, \mathrm{high\_speed}\}$ is the discretisation of the row by $s_t$. The contextual gate is the conjunction
\begin{equation}
G_{\mathrm{context}}(x_t) = G_{\mathrm{recent\text{-}drop}}(x_t) \;\wedge\; \neg\, G_{\mathrm{refuel}}(x_t) \;\wedge\; \neg\, G_{\mathrm{slosh}}(x_t)
\label{eq:gate_conjunction}
\end{equation}
with the three constituent predicates defined as
\begin{align}
G_{\mathrm{recent\text{-}drop}}(x_t) &= \mathbb{1}\!\left[\, \frac{1}{12} \sum_{i = t-11}^{t} \max\bigl(0,\, -\Delta F_i\bigr) \geq 1.0\% \,\right] \label{eq:gate_recent_drop}\\
G_{\mathrm{refuel}}(x_t) &= \mathbb{1}\bigl[\, (t - t_{\mathrm{lastRefuel}}) \leq 300\text{ s} \,\bigr], \quad t_{\mathrm{lastRefuel}} := t'\ \text{such that}\ \Delta F_{t'} \geq +5\% \label{eq:gate_refuel}\\
G_{\mathrm{slosh}}(x_t) &= \mathbb{1}\!\left[\, \frac{L_t / L_{\max}}{s_t / 80 + 0.1} > 5 \,\right] \label{eq:gate_slosh}
\end{align}
where $\Delta F_i$ is the fuel-delta feature at row $i$, $t_{\mathrm{lastRefuel}}$ is the sample index of the most recent $+5\%$ fuel jump on the current trip, and $L_{\max} = 100$ is the load normalisation. Equation~\ref{eq:gate_recent_drop} requires that a real recent drop is in progress before an alert is admitted (this suppresses spurious point-anomalies without underlying fuel movement). Equation~\ref{eq:gate_refuel} exempts the five-minute post-refuel quiescence window during which the sender is settling. Equation~\ref{eq:gate_slosh} suppresses the specific high-load-at-low-speed regime (uphill grind, cargo overload) that produced repeatable false positives in the 2026-06-23 TRCK-03 trip and motivated the addition of \texttt{load\_per\_speed\_norm} as feature~23.

\subsection{Hybrid Fusion Methodology}
Two fusion strategies are characterised in the inference pipeline. The Hybrid Union (Isolation Forest OR Matrix Profile) maximises recall. The Hybrid Intersection used for production alerting is a trip-level Matrix Profile gate over row-level Isolation Forest predictions: a row is flagged if and only if (a) the Isolation Forest flagged that row and (b) the trip's overall Matrix Profile decision was ``anomalous shape detected''. The trip-gated definition preserves the precision intent of the intersection while avoiding the granularity-mismatch penalty of a strict row-level logical AND.

\subsection{Contextual Gate Methodology}
The contextual gate in the backend receives every row flagged by the ML layer and applies three deterministic rules: drop (suppress repeated flags inside the same window), drive-rise (suppress flags during legitimate post-acceleration recovery), and parked-refuel (suppress flags during the 5-minute post-refuel quiescence window). Rows that pass the gate are promoted into user-facing fuel-anomaly alerts, while suppressed rows remain available as model-layer flags for audit and debugging.

\subsection{Backend Inference Integration}
The ML service is a Flask microservice exposing a detection endpoint. After each telemetry insert, the backend asynchronously invokes anomaly detection; the result is written back to \texttt{telemetry\_logs} as \texttt{anomaly\_flag}, \texttt{anomaly\_score}, and \texttt{model\_\allowbreak source}, and a row is appended to the \texttt{alerts} table with \texttt{alert\_type = 'fuel\_\allowbreak anomaly'} when the contextual gate permits.

\section{Evaluation Methodology}
\label{sec:evaluation}

This section defines the evaluation design used to validate the system against its specific objectives. It describes the cross-validation protocol applied to the anomaly-detection models, the mapping between each objective and its evaluation protocol and metric, and the formal definitions of the performance metrics reported in Chapter~6.

\subsection{Cross-Validation Design}

All Chapter 6 SO4 numbers are produced by a leave-one-trip-out (LOTO) evaluation that holds the calibration set strictly disjoint from the test set.\glsadd{not:Dtrain}\glsadd{not:Dtest} For each of the twelve project-native trips, the three Isolation Forest variants and the Matrix Profile detector are refit and recalibrated on the clean rows of the remaining eleven project-native trips only; the held-out trip is then scored with that fold's models. Each prediction is therefore strictly out-of-sample with respect to both model fitting and threshold calibration. The Isolation Forest threshold for every variant is the 0.985 quantile\glsadd{not:q0985} of the fold's clean training scores, internally consistent with the contamination = 0.015 hyperparameter; no precision margin and no per-context shimming are applied.

\subsection{Objective-Based Evaluation Protocol}

Table~\ref{tab:objective_protocol_metric} maps each specific objective to the evaluation protocol used to test it and the metrics reported for it.

\begin{table}[!htbp]
\centering
\caption{Objective, Protocol, and Metric}
\label{tab:objective_protocol_metric}
\scriptsize
\begin{tabular}{|p{0.18\textwidth}|p{0.48\textwidth}|p{0.24\textwidth}|}
\hline
\textbf{Objective} & \textbf{Evaluation protocol} & \textbf{Metric} \\
\hline
SO1 - Fuel accuracy\glsadd{not:Vsensori}\glsadd{not:Vmanuali}\glsadd{not:Ctank}\glsadd{not:Vpumpi}\glsadd{not:VHMIi}\glsadd{not:ei} & Known-volume pump-gas events + manufacturer dashboard supporting readings + per-tank-band drift/noise/jitter on $n=20,363$ valid samples & MAPE, max error, mean absolute noise (\%), drift at speed = 0 \\
\hline
SO2 - Communication & LoRa distance test (10 m to 5 km), GSM/LTE drive test at five locations, three-region packet-delivery drive (urban/highway/dead-zone) & Median + p95 latency, packet-success fraction, SPIFFS-replay share, RSSI/SNR per band \\
\hline
SO3 - GPS & \texttt{sent\_at} interval analysis on field-test telemetry; GPS fix availability; route vs OSRM\glsadd{OSRM} polyline comparison & Median + p95 inter-sample dt, fix availability (\%), per-segment fuel (\% per km) \\
\hline
SO4 - ML & LOTO cross-validation on \texttt{project\_native\_\allowbreak injected} + threshold/IQR baseline + Hybrid Union/Intersection ablation + OOD stress test on \texttt{public\_clean} + extended on-road validation & Precision, recall, FPR, F1, Wilson 95\% CI, per-model confusion matrices, per-class recall \\
\hline
SO5 - Alerts & Alert-trigger verification across controlled-injection and field-test trip sets & Trigger correctness (TP rate), driver-response capture, HMI/dashboard/mobile sync correctness \\
\hline
\end{tabular}
\normalsize
\end{table}

Table~\ref{tab:objective_protocol_metric} is the evaluation contract for the thesis. It states what evidence is accepted for each objective before the results are presented, which prevents the results chapter from choosing metrics after the fact. The same mapping is used in Chapter~6: SO1 is judged by fuel error and noise, SO2 by latency and packet delivery, SO3 by GPS cadence and fix availability, SO4 by classifier metrics, and SO5 by alert correctness and interface synchronization.

\subsection{Evaluation Metrics}
\label{sec:evaluate}

To evaluate the anomaly-detection models, predicted labels were compared with the evaluation-corpus labels to count true positives (TP), false positives (FP), true negatives (TN), and false negatives (FN). The key metrics were computed as follows:

\textbf{Precision}
\begin{equation}
\mathrm{Precision} = \frac{\gls{not:TP}}{\gls{not:TP} + \gls{not:FP}}
\label{eq:precision}
\end{equation}

Precision reflects how many detected anomalies are truly anomalous. A high precision means fewer false alarms, which is important for maintaining trust in the alerting system.

\textbf{False Positive Rate (FPR)}
\begin{equation}
\mathrm{FPR} = \frac{\gls{not:FP}}{\gls{not:FP} + \gls{not:TN}}
\label{eq:fpr}
\end{equation}

FPR quantifies the likelihood of incorrectly classifying normal behavior as anomalous. Lower FPR values are desired to reduce noise in the system's alerts.

\textbf{Recall}
\begin{equation}
\mathrm{Recall} = \frac{\gls{not:TP}}{\gls{not:TP} + \gls{not:FN}}
\label{eq:recall}
\end{equation}

Recall reflects how many true anomalies were detected by the model.

\textbf{F1 Score}
\begin{equation}
\mathrm{F1} = \frac{2 \cdot \mathrm{Precision} \cdot \mathrm{Recall}}{\mathrm{Precision} + \mathrm{Recall}}
\label{eq:f1}
\end{equation}

F1 score provides a single balance between precision and recall, which is useful when anomalous rows are rare.

\textbf{Accuracy}
\begin{equation}
\mathrm{Accuracy} = \frac{\gls{not:TP} + \gls{not:TN}}{\gls{not:TP} + \gls{not:TN} + \gls{not:FP} + \gls{not:FN}}
\label{eq:accuracy}
\end{equation}

Accuracy provides an overall view of model correctness but can be misleading in imbalanced datasets, which is why it's used alongside precision and FPR.

All these metrics use the following classification outcomes:
\begin{align*}
\text{TP (True Positives)} &= \text{Number of correctly identified anomalies}, \\
\text{FP (False Positives)} &= \text{Number of normal events incorrectly flagged as anomalies}, \\
\text{TN (True Negatives)} &= \text{Number of correctly identified normal cases}, \\
\text{FN (False Negatives)} &= \text{Anomalies that were missed by the model}.
\end{align*}

\textbf{Detection Latency}

\begin{equation}
\gls{not:L}_{\mathrm{det}} = \gls{not:Tdetected} - \gls{not:Tactual}
\label{eq:latency-detection}
\end{equation}

where:
\begin{align*}
T_{\text{detected}} &= \text{the time the anomaly was identified by the system}, \\
T_{\text{actual}} &= \text{the time the anomaly truly occurred}.
\end{align*}

The SO4 objective is achieved if the unsupervised system attains at least 90\% precision, at least 85\% recall, and a false positive rate not exceeding 10\%. Detection latency is evaluated separately to characterize how quickly the backend and alerting surfaces respond after an anomalous event occurs.

\section{Data Collection and Research-Ethics Procedure}

The data-collection procedure followed a non-invasive deployment model. The prototype was attached through the vehicle OBD-II port and external HMI wiring, and it did not require fuel-tank opening, sender-loom modification, or direct control of vehicle operation. Controlled fuel-anomaly labels used for model evaluation were produced through the methodology described in the machine-learning section rather than by asking drivers to drain fuel from operational fleet vehicles.

\subsection{Study Procedures}

Upon voluntary agreement to participate, the study procedures were structured as follows:

\begin{enumerate}
    \item Before a monitored trip, a member of the research team installed the monitoring device or IoT module in the fleet vehicle assigned to the participant.
    \item Once installed, the device automatically recorded OBD-II fuel readings, distance, travel time, GPS location, communication channel, and alert status during normal vehicle operation.
    \item Participants were not required to change their normal driving behavior during the data-collection period.
    \item After the data-collection period, a member of the research team removed the monitoring device from the vehicle.
    \item The study collected fleet-administrator feedback on device functionality and operational issues. The survey was not administered to drivers and did not ask personal questions directed at them.
\end{enumerate}

\subsection{Risks and Discomforts}

This study involves minimal risk. The following potential risks and discomforts have been identified:

\begin{itemize}
    \item The monitoring device will not control, modify, or interfere with the vehicle's systems and will not affect driving performance or safety. It will passively record only fuel-level data, GPS data, and travel information.
    \item Participants may experience minor inconvenience or brief disruption during the installation and removal of the monitoring device before and after the trip.
    \item GPS location data will be collected during normal operations, which may raise privacy concerns. However, the data will be anonymized through a unique code number assigned to each dataset and will be shared only with authorized members of the research team and Hino Paranaque's fleet administrator for academic research purposes.
    \item No foreseeable psychological, physical, or occupational risks are expected to result from participation in this study.
\end{itemize}

\subsection{Steps to Minimize Risks}

The research team undertook the following measures to minimize and mitigate potential risks:

\begin{itemize}
    \item The monitoring device will be installed and removed only by a member of the research team to ensure that the vehicle's systems are not disrupted.
    \item All collected data will be anonymized by replacing participant names with unique code numbers before analysis.
    \item Access to raw data will be restricted to authorized members of the research team and Hino Paranaque's fleet administrator only.
    \item All digital data will be stored in a secure, password-protected system to prevent unauthorized access.
    \item Physical documents, such as signed consent forms, will be kept in a locked cabinet accessible only to the principal investigator and faculty adviser.
    \item Data will be used only for academic purposes, and only summarized and anonymized results will be included in reports or publications.
\end{itemize}

\subsection{Benefits}

Participants will not receive direct financial benefits from participation. However, the results of this study may contribute to the following non-monetary benefits:

\begin{itemize}
    \item \textbf{For the partner company and fleet managers:} The system may improve fuel efficiency monitoring and support the detection and prevention of fuel misuse, leakage, or theft.
    \item \textbf{For the transportation and logistics industry:} The study contributes to research on non-parametric machine learning approaches for fleet management, which may inform future industry practices.
    \item \textbf{For the public and the environment:} Improved fuel monitoring systems may help reduce fuel waste and contribute to lower carbon emissions.
    \item \textbf{For the research community:} The findings may add to academic knowledge on IoT-based monitoring systems and anomaly detection methods in real-world applications.
\end{itemize}

\subsection{Compensation}

Participants will not receive any compensation, tokens, or gifts for their participation in this study.

\subsection{Data Storage, Duration, and Disposal}

\textbf{Storage of Digital Data:} All digital data collected by the monitoring device will be stored in a Supabase cloud database, a secure and password-protected platform accessible only to authorized members of the research team. Files and database entries will be protected by unique credentials known only to the principal investigator and faculty adviser.

\textbf{Storage of Physical Data:} Signed consent forms and any printed documents will be stored in a locked filing cabinet located at the research team's office. Only the principal investigator and faculty adviser will have access to these physical documents.

\textbf{Duration of Storage:} All digital and physical data will be retained for five (5) years from the date of completion of the study, in accordance with De La Salle University's data retention policy.

\textbf{Disposal:} After the retention period, all digital data will be permanently deleted from the Supabase database and any local copies using secure deletion procedures. All physical documents will be cross-shredded and disposed of in a secure manner.

\subsection{Confidentiality}

\textbf{For Digital Data:} Participant names will not appear in any reports, datasets, publications, or presentations. A unique code number will be assigned to each participant's data in place of the participant's name. All digital files will be stored in a password-protected Supabase cloud database accessible only to the research team. Data will be used strictly for academic research purposes, and only summarized and anonymized results will be presented in the final thesis or any related publications.

\textbf{For Physical Data:} Signed consent forms and any printed documents will be stored in a locked cabinet. Only the principal investigator and faculty adviser will have access to these hard copy documents.

\subsection{Withdrawal Procedure}

Participation in this study is entirely voluntary, and participants may withdraw at any time without penalty. To withdraw, participants may follow these steps:

\begin{enumerate}
    \item Inform the principal investigator, Samuelle P. Barja, of the decision to withdraw, either verbally or in writing.
    \item Contact the principal investigator through the following:
    \begin{itemize}
        \item samuelle\_barja@dlsu.edu.ph
        \item 09564054471
    \end{itemize}
    \item Upon notification of withdrawal, the research team will immediately stop collecting data from the participant's assigned vehicle.
    \item Any data already collected that is identifiable to the participant will be permanently deleted from all storage systems within a reasonable timeframe.
    \item Participants will not be required to provide any reason for withdrawal, and the decision to withdraw will not result in any consequence to their employment or relationship with their employer.
\end{enumerate}

\section{Summary}

This chapter presented the implementation and evaluation methodology for the IoT-enabled fuel monitoring, GPS tracking, and anomaly detection prototype. The approach integrates OBD-II fuel acquisition, edge filtering, prioritized communication with offline replay, Supabase-backed state synchronization, Isolation Forest and Matrix Profile anomaly detection, objective-based evaluation metrics, and research-ethics controls for field data collection.
